% Generated by tools/build.py -- do not edit; edit data/data-structures-algorithms/*.txt
% No preamble here: this file is \input by main.tex.
\chapter{Unit 7: Data Structures and Algorithms}

\section{Arrays, Sparse Matrices, Stacks and Queues}

\begin{question}{H}
An array A[1..10][1..15] of 2-byte elements is stored in row-major order starting at address 100. The address of A[5][7] is
\optsshort{232}{240}{228}{234}
\end{question}

\begin{question}{H}
An array A[1..10][1..15] of 2-byte elements is stored in column-major order starting at address 100. The address of A[5][7] is
\optsshort{236}{228}{244}{232}
\end{question}

\begin{question}{H}
An array A[$-$5..5][2..10] of 4-byte elements is stored in row-major order starting at address 1000. The address of A[0][5] is
\optsshort{1180}{1216}{1192}{1200}
\end{question}

\begin{question}{M}
A lower-triangular 10 $\times$ 10 matrix is stored row by row, keeping only the elements on and below the diagonal. The number of elements stored is
\optsshort{50}{45}{55}{100}
\end{question}

\begin{question}{M}
Elements on and below the diagonal of a lower-triangular matrix A[1..n][1..n] are stored row by row in a one-dimensional array starting at index 0. The index of A[5][3] is
\optsshort{13}{15}{10}{12}
\end{question}

\begin{question}{M}
In the triplet (row, column, value) representation of a sparse matrix with t non-zero elements, including one header triplet for the dimensions, the number of values stored is
\optsshort{3t}{t + 3}{2t + 1}{3(t + 1)}
\end{question}

\begin{question}{H}
The value of the postfix expression 6 2 3 + $-$ 3 8 2 / + * 2 ^ 3 + is
\optsshort{55}{52}{64}{49}
\end{question}

\begin{question}{H}
The value of the postfix expression 5 6 2 + * 12 4 / $-$ is
\optsshort{43}{37}{46}{40}
\end{question}

\begin{question}{H}
The postfix form of the infix expression A + B * C $-$ D / E is
\optsshort{ABC+*DE/$-$}{AB+C*DE/$-$}{ABC*+DE$-$/}{ABC*+DE/$-$}
\end{question}

\begin{question}{H}
The postfix form of A + B * (C ^ D $-$ E) ^ (F + G * H) $-$ I is
\optsshort{ABCDE\textsuperscript{$-$FGH}*+^*+I$-$}{ABCD\textsuperscript{E}$-$FGH*+^+*I$-$}{ABCD\textsuperscript{E}$-$FGH*+^*+I$-$}{AB+CD\textsuperscript{E}$-$*FGH*+\textsuperscript{I}$-$}
\end{question}

\begin{question}{H}
The prefix form of the infix expression A + B * C $-$ D / E is
\optsshort{$-$+A*BC/DE}{+$-$A*BC/DE}{+A$-$*BC/DE}{$-$+*ABC/DE}
\end{question}

\begin{question}{M}
Since ^ is right-associative, the postfix form of A ^ B ^ C is
\optsshort{ABC^ ^ (with A\textsuperscript{B} first)}{ABC^^}{AB\textsuperscript{C}^}{A\textsuperscript{BC}^}
\end{question}

\begin{question}{M}
The postfix form of (A $-$ B) / C * D + E is
\optsshort{AB$-$CD*/E+}{AB$-$C/D*E+}{ABC$-$/D*E+}{AB$-$C/DE*+}
\end{question}

\begin{question}{M}
During the evaluation of the postfix expression 2 3 4 * + 5 $-$ using a stack, the maximum number of elements on the stack at any time is
\optsshort{3}{2}{5}{4}
\end{question}

\begin{question}{H}
The numbers 1, 2, 3, 4, 5 are pushed onto a stack in that order, and pops may occur at any time. Which output sequence is NOT possible?
\optsshort{5, 4, 3, 2, 1}{4, 5, 3, 1, 2}{3, 2, 4, 1, 5}{2, 1, 5, 4, 3}
\end{question}

\begin{question}{H}
The number of different output permutations obtainable by passing 1, 2, 3, 4 through a stack (pushing in order, popping at any time) is
\optsshort{42}{14}{24}{16}
\end{question}

\begin{question}{H}
Given below are two statements.\\
Statement I: A stack is used to implement function calls and recursion.\\
Statement II: A queue is used for breadth-first search of a graph.
\optslong{Both Statement I and Statement II are true}{Both Statement I and Statement II are false}{Statement I is true but Statement II is false}{Statement I is false but Statement II is true}
\end{question}

\begin{question}{H}
In a circular queue of size n implemented with front and rear indices, where one slot is kept empty, the queue is full when
\optsshort{(rear + 1) \% n == front}{(front + 1) \% n == rear}{rear == front}{rear == n $-$ 1}
\end{question}

\begin{question}{M}
A queue is implemented using two stacks, where dequeue moves all elements from the input stack to the output stack only when the output stack is empty. The amortized cost per operation is
\optsshort{O(1)}{O(n)}{O(n\textsuperscript{2})}{O(log n)}
\end{question}

\begin{question}{M}
Match List I with List II.\\
List I (Data structure): A. Stack  B. Queue  C. Deque  D. Priority queue\\
List II (Removal rule): I. From either end  II. Element with the highest priority  III. Last in, first out  IV. First in, first out
\optsshort{A-III, B-IV, C-I, D-II}{A-IV, B-III, C-I, D-II}{A-III, B-IV, C-II, D-I}{A-III, B-I, C-IV, D-II}
\end{question}

\begin{question}{M}
An input-restricted deque allows
\optslong{insertion and deletion at one end only}{no deletions}{insertion at both ends but deletion at one end}{insertion at one end but deletion at both ends}
\end{question}

\begin{question}{H}
The minimum number of moves needed to solve the Tower of Hanoi with 6 disks is
\optsshort{63}{36}{64}{127}
\end{question}

\begin{question}{M}
Which data structure is most suitable for checking whether the parentheses in an expression are balanced?
\optsshort{Queue}{Stack}{Binary search tree}{Hash table}
\end{question}

\begin{question}{M}
Given below are two statements: one is labelled as Assertion A and the other is labelled as Reason R.\\
Assertion A: Accessing the element at index i of an array takes constant time.\\
Reason R: An array stores its elements in contiguous memory, so the address can be computed as base + i $\times$ size.
\optslong{Both A and R are true and R is the correct explanation of A}{Both A and R are true but R is NOT the correct explanation of A}{A is true but R is false}{A is false but R is true}
\end{question}

\begin{question}{M}
Inserting an element at the beginning of an array of n elements requires, in the worst case,
\optsshort{O(1) time}{O(n\textsuperscript{2}) time}{O(log n) time}{O(n) time}
\end{question}

\begin{question}{M}
A stack implemented with an array of size n overflows when
\optslong{top == 0}{the stack is empty and a pop is attempted}{top == n $-$ 1 and a push is attempted}{top == $-$1}
\end{question}

\begin{question}{M}
Given below are two statements.\\
Statement I: A queue can be implemented using two stacks.\\
Statement II: A stack can be implemented using two queues.
\optslong{Both Statement I and Statement II are true}{Both Statement I and Statement II are false}{Statement I is true but Statement II is false}{Statement I is false but Statement II is true}
\end{question}

\begin{question}{M}
A queue implemented as a singly linked list with pointers to the front and rear nodes supports
\optslong{enqueue and dequeue in O(1)}{enqueue in O(n) and dequeue in O(1)}{enqueue in O(1) and dequeue in O(n)}{both in O(n)}
\end{question}

\begin{question}{M}
Which of the following applications does NOT normally use a stack?
\optslong{CPU scheduling with round robin}{Undo operation in an editor}{Evaluating postfix expressions}{Backtracking}
\end{question}

\begin{question}{M}
Given below are two statements.\\
Statement I: Converting an infix expression to postfix requires a queue of operators.\\
Statement II: A postfix expression needs parentheses to specify the order of evaluation.
\optslong{Both Statement I and Statement II are true}{Both Statement I and Statement II are false}{Statement I is true but Statement II is false}{Statement I is false but Statement II is true}
\end{question}

\begin{question}{M}
The number of push operations needed to evaluate the postfix expression 2 3 + 4 5 * $-$ (counting the pushes of results as well) is
\optsshort{4}{5}{9}{7}
\end{question}

\begin{question}{M}
Match List I with List II.\\
List I (Operation on an array of n elements): A. Access by index  B. Linear search  C. Binary search (sorted)  D. Insertion at a given position\\
List II (Worst-case time): I. O(log n)  II. O(n) due to shifting  III. O(1)  IV. O(n)
\optsshort{A-III, B-IV, C-II, D-I}{A-III, B-IV, C-I, D-II}{A-III, B-I, C-IV, D-II}{A-IV, B-III, C-I, D-II}
\end{question}

\begin{question}{M}
A recursive function calls itself n times before reaching the base case. The space used on the call stack is
\optsshort{O(n\textsuperscript{2})}{O(1)}{O(log n)}{O(n)}
\end{question}

\begin{question}{M}
Given below are two statements: one is labelled as Assertion A and the other is labelled as Reason R.\\
Assertion A: In a circular queue, the space freed by dequeued elements can be reused.\\
Reason R: A queue follows the first-in, first-out principle.
\optslong{Both A and R are true and R is the correct explanation of A}{Both A and R are true but R is NOT the correct explanation of A}{A is true but R is false}{A is false but R is true}
\end{question}

\begin{question}{M}
Given below are two statements.\\
Statement I: A sparse matrix is one in which most of the elements are non-zero.\\
Statement II: Storing a sparse matrix as a linked list or as triplets saves memory compared with a two-dimensional array.
\optslong{Both Statement I and Statement II are true}{Both Statement I and Statement II are false}{Statement I is true but Statement II is false}{Statement I is false but Statement II is true}
\end{question}

\begin{question}{M}
Which of the following expressions is in prefix (Polish) notation?
\optsshort{AB+C*}{+*ABC}{A B C + *}{A+B*C}
\end{question}

\begin{question}{M}
The value of the prefix expression $-$ + 4 * 2 3 / 8 4 is
\optsshort{8}{12}{10}{14}
\end{question}

\begin{question}{L}
Which data structure is used by the operating system for buffering keystrokes and print jobs?
\optsshort{Stack}{Tree}{Queue}{Graph}
\end{question}

\begin{question}{M}
A double-ended queue (deque) supports insertion and deletion at both ends. Which of the following can it simulate directly?
\optsshort{Only a stack}{Only a queue}{Both a stack and a queue}{Neither a stack nor a queue}
\end{question}

\begin{question}{M}
Given below are two statements.\\
Statement I: FORTRAN stores arrays in row-major order.\\
Statement II: In row-major order, the elements of each row are stored contiguously.
\optslong{Both Statement I and Statement II are true}{Both Statement I and Statement II are false}{Statement I is true but Statement II is false}{Statement I is false but Statement II is true}
\end{question}

\begin{answers}
\ans{1}{1}{100 + 2 $\times$ ((5 $-$ 1) $\times$ 15 + (7 $-$ 1)) = 100 + 2 $\times$ 66 = 232.}
\ans{2}{2}{100 + 2 $\times$ ((7 $-$ 1) $\times$ 10 + (5 $-$ 1)) = 100 + 2 $\times$ 64 = 228.}
\ans{3}{3}{Each row has 9 elements: 1000 + 4 $\times$ ((0 $-$ ($-$5)) $\times$ 9 + (5 $-$ 2)) = 1000 + 4 $\times$ 48 = 1192.}
\ans{4}{3}{n(n + 1)/2 = 10 $\times$ 11 / 2 = 55.}
\ans{5}{4}{Rows 1 to 4 hold 1 + 2 + 3 + 4 = 10 elements, so A[5][3] is at 10 + (3 $-$ 1) = 12.}
\ans{6}{4}{The header stores the number of rows, columns and non-zero elements.}
\ans{7}{2}{6 $-$ 5 = 1; 8/2 + 3 = 7; 1 $\times$ 7 = 7; 7\textsuperscript{2} = 49; 49 + 3 = 52.}
\ans{8}{2}{5 $\times$ (6 + 2) = 40; 12/4 = 3; 40 $-$ 3 = 37.}
\ans{9}{4}{* and / are applied before + and $-$, and operators of equal precedence associate left to right.}
\ans{10}{3}{^ has the highest precedence and is right-associative.}
\ans{11}{1}{The last operator applied is $-$, whose operands are A + B * C and D / E.}
\ans{12}{2}{A ^ B ^ C means A ^ (B ^ C), giving ABC^^.}
\ans{13}{2}{/ and * are left-associative: ((A $-$ B) / C) * D, then + E.}
\ans{14}{1}{2, 3 and 4 are pushed before the first operator is applied.}
\ans{15}{2}{After 4, 5 and 3 are popped, 2 is above 1 on the stack, so 1 cannot come out before 2.}
\ans{16}{2}{It is the Catalan number C(4) = (1/5) $\times$ C(8, 4) = 14.}
\ans{17}{1}{Depth-first search, in contrast, uses a stack (explicitly or through recursion).}
\ans{18}{1}{Keeping one slot empty distinguishes ``full'' from ``empty'' (front == rear); the queue holds at most n $-$ 1 elements.}
\ans{19}{1}{Each element is pushed and popped at most twice in total.}
\ans{20}{1}{Basic linear data structures.}
\ans{21}{4}{An output-restricted deque allows deletion at one end and insertion at both ends.}
\ans{22}{1}{2\textsuperscript{n} $-$ 1 = 64 $-$ 1 = 63.}
\ans{23}{2}{Each opening bracket is pushed and matched with the next closing bracket.}
\ans{24}{1}{This direct address computation is what gives arrays random access.}
\ans{25}{4}{All n elements must be shifted one position.}
\ans{26}{3}{Popping from an empty stack (top == $-$1) is underflow.}
\ans{27}{1}{Both are possible, with different costs for push and pop.}
\ans{28}{1}{Insert at the rear and delete at the front.}
\ans{29}{1}{Round-robin scheduling uses a circular queue.}
\ans{30}{2}{Infix-to-postfix conversion uses a stack of operators, and postfix expressions need no parentheses.}
\ans{31}{4}{Four operands are pushed, plus three results (for +, * and $-$): 7 pushes.}
\ans{32}{2}{Basic array operation costs.}
\ans{33}{4}{Each pending call keeps its own activation record.}
\ans{34}{2}{Both are true, but reuse comes from the indices wrapping around modulo the array size.}
\ans{35}{4}{A sparse matrix has mostly zero elements, which is why compact representations save space.}
\ans{36}{2}{In prefix notation the operator precedes its operands.}
\ans{37}{1}{* 2 3 = 6; + 4 6 = 10; / 8 4 = 2; $-$ 10 2 = 8.}
\ans{38}{3}{Requests are served in arrival order.}
\ans{39}{3}{Using one end only gives a stack; inserting at one end and deleting at the other gives a queue.}
\ans{40}{4}{FORTRAN uses column-major order; C uses row-major order.}
\end{answers}

\section{Linked Lists, Priority Queues and Heaps}

\begin{question}{H}
Match List I with List II.\\
List I (Operation on a singly linked list of n nodes with only a head pointer): A. Insert at the beginning  B. Search for a key  C. Access the kth node  D. Sort the list with merge sort\\
List II (Time): I. O(1)  II. O(n)  III. O(k)  IV. O(n log n)
\optsshort{A-I, B-III, C-II, D-IV}{A-I, B-II, C-IV, D-III}{A-I, B-II, C-III, D-IV}{A-II, B-I, C-III, D-IV}
\end{question}

\begin{question}{H}
In a doubly linked list, given a pointer p to a node to be deleted (not the first or last node), which statements delete it?
\optslong{p->prev = p->next; p->next = p->prev;}{p->prev->next = p->next; p->next->prev = p->prev;}{p->next->next = p->prev; p->prev->prev = p->next;}{p->next = NULL; p->prev = NULL;}
\end{question}

\begin{question}{H}
Given below are two statements.\\
Statement I: Binary search can be applied efficiently to a sorted singly linked list.\\
Statement II: Finding the middle element of a singly linked list requires random access.
\optslong{Both Statement I and Statement II are true}{Both Statement I and Statement II are false}{Statement I is true but Statement II is false}{Statement I is false but Statement II is true}
\end{question}

\begin{question}{H}
Floyd's cycle-finding algorithm detects a loop in a linked list using
\optslong{sorting the nodes}{a hash table of visited nodes}{two pointers moving at different speeds}{recursion only}
\end{question}

\begin{question}{M}
Reversing a singly linked list of n nodes in place requires
\optslong{O(n) time and O(1) extra space}{O(log n) time}{O(n\textsuperscript{2}) time}{O(1) time}
\end{question}

\begin{question}{M}
In a circular singly linked list with a pointer only to the last node, which operations take O(1) time?\\
A. Insert at the beginning  B. Insert at the end  C. Delete the first node  D. Delete the last node\\
Choose the correct answer from the options given below:
\optsshort{A, B and C only}{A, B, C and D}{C and D only}{A and B only}
\end{question}

\begin{question}{M}
What does the following function print for a linked list 1 $\rightarrow$ 2 $\rightarrow$ 3 $\rightarrow$ 4?\\
void f(Node *h) \{ if (h == NULL) return; f(h->next); printf(``\%d '', h->data); \}
\optsshort{1 2 3 4}{4 2}{1 3}{4 3 2 1}
\end{question}

\begin{question}{M}
An XOR linked list stores in each node
\optslong{the XOR of the addresses of the previous and next nodes}{two pointers}{the address of the next node only}{the XOR of the data values}
\end{question}

\begin{question}{M}
Given below are two statements: one is labelled as Assertion A and the other is labelled as Reason R.\\
Assertion A: A linked list uses memory more flexibly than an array when the number of elements varies.\\
Reason R: Linked list nodes must be stored in consecutive memory locations.
\optslong{Both A and R are true and R is the correct explanation of A}{Both A and R are true but R is NOT the correct explanation of A}{A is true but R is false}{A is false but R is true}
\end{question}

\begin{question}{H}
In a max-heap stored in a 1-indexed array, the parent and children of the node at index i are at
\optsshort{i/2, 2i and 2i + 1}{2i, i/2 and i/2 + 1}{i $-$ 1, i + 1 and i + 2}{i/2, i and i + 1}
\end{question}

\begin{question}{H}
Building a max-heap from the array [4, 1, 3, 2, 16, 9, 10, 14, 8, 7] using the bottom-up BUILD-MAX-HEAP procedure gives
\optslong{[16, 10, 14, 8, 7, 9, 3, 2, 4, 1]}{[16, 14, 10, 4, 7, 9, 3, 2, 8, 1]}{[16, 14, 10, 8, 7, 9, 3, 2, 4, 1]}{[16, 14, 9, 8, 7, 10, 3, 2, 4, 1]}
\end{question}

\begin{question}{H}
From the max-heap [16, 14, 10, 8, 7, 9, 3, 2, 4, 1], the maximum is deleted (the last element is moved to the root and sifted down). The resulting heap is
\optslong{[14, 8, 10, 2, 7, 9, 3, 4, 1]}{[10, 14, 8, 4, 7, 9, 3, 2, 1]}{[14, 10, 8, 7, 4, 9, 3, 2, 1]}{[14, 8, 10, 4, 7, 9, 3, 2, 1]}
\end{question}

\begin{question}{H}
The keys 10, 20, 15, 30, 40 are inserted one by one into an initially empty max-heap. The resulting array is
\optsshort{[40, 30, 20, 15, 10]}{[40, 30, 15, 20, 10]}{[40, 30, 15, 10, 20]}{[40, 20, 30, 10, 15]}
\end{question}

\begin{question}{H}
The time complexity of building a binary heap from n arbitrary elements using the bottom-up method is
\optsshort{O(n)}{O(n log n)}{O(n\textsuperscript{2})}{O(log n)}
\end{question}

\begin{question}{H}
Match List I with List II.\\
List I (Operation on a binary max-heap): A. Find maximum  B. Insert one element  C. Build a heap from n elements  D. Heap sort of n elements\\
List II (Worst-case time): I. O(n log n)  II. O(log n)  III. O(1)  IV. O(n)
\optsshort{A-II, B-III, C-IV, D-I}{A-III, B-IV, C-II, D-I}{A-III, B-II, C-I, D-IV}{A-III, B-II, C-IV, D-I}
\end{question}

\begin{question}{M}
In a max-heap of n distinct elements, the minimum element is found in
\optslong{the last level only at index n}{one of the leaves, which requires examining about n/2 elements}{the root}{the left child of the root}
\end{question}

\begin{question}{M}
The number of leaves in a binary heap with n elements is
\optsshort{$\lfloor$n/2$\rfloor$}{n $-$ 1}{$\lceil$n/2$\rceil$}{log2 n}
\end{question}

\begin{question}{H}
The number of distinct max-heaps that can be built from the keys 1, 2, 3, 4, 5 is
\optsshort{16}{5}{8}{20}
\end{question}

\begin{question}{M}
The number of distinct max-heaps that can be built from 7 distinct keys is
\optsshort{48}{120}{80}{42}
\end{question}

\begin{question}{M}
The height of a binary heap with n elements is
\optsshort{n/2}{log2 n squared}{$\surd$n}{$\lfloor$log2 n$\rfloor$}
\end{question}

\begin{question}{M}
Given below are two statements.\\
Statement I: Heap sort is a stable sorting algorithm.\\
Statement II: Heap sort runs in O(n log n) time in the worst case and sorts in place.
\optslong{Both Statement I and Statement II are true}{Both Statement I and Statement II are false}{Statement I is true but Statement II is false}{Statement I is false but Statement II is true}
\end{question}

\begin{question}{M}
k sorted lists containing n elements in total are merged into one sorted list using a min-heap of size k. The time required is
\optsshort{O(n log k)}{O(k log n)}{O(n)}{O(nk)}
\end{question}

\begin{question}{M}
The kth smallest element in a binary min-heap of n elements can be found, using an auxiliary heap, in
\optsshort{O(k\textsuperscript{2})}{O(k log k)}{O(n)}{O(n log n)}
\end{question}

\begin{question}{M}
Match List I with List II.\\
List I (Heap variant): A. Binomial heap  B. Fibonacci heap  C. d-ary heap\\
List II (Feature): I. Height O(log\_d n), cheaper decrease-key and more expensive delete-min  II. O(1) amortized decrease-key, used to speed up Dijkstra and Prim  III. Collection of binomial trees supporting union in O(log n)
\optsshort{A-III, B-II, C-I}{A-II, B-III, C-I}{A-I, B-II, C-III}{A-III, B-I, C-II}
\end{question}

\begin{question}{M}
Given below are two statements: one is labelled as Assertion A and the other is labelled as Reason R.\\
Assertion A: A binary heap is usually stored in an array without pointers.\\
Reason R: A heap is a complete binary tree, so the positions of the parent and children can be computed from the index.
\optslong{Both A and R are true and R is the correct explanation of A}{Both A and R are true but R is NOT the correct explanation of A}{A is true but R is false}{A is false but R is true}
\end{question}

\begin{question}{M}
In a min-heap stored in an array, which of the following arrays is a valid min-heap?
\optsshort{[2, 1, 3, 4, 5, 6]}{[1, 5, 2, 4, 7, 3]}{[1, 2, 3, 0, 4, 5]}{[1, 3, 2, 7, 4, 5]}
\end{question}

\begin{question}{M}
Which of the following is a valid max-heap?
\optsshort{[25, 14, 13, 16, 10, 8, 12]}{[25, 14, 16, 13, 10, 8, 12]}{[25, 14, 12, 13, 10, 8, 16]}{[25, 12, 16, 13, 10, 8, 14]}
\end{question}

\begin{question}{M}
A priority queue is most efficiently implemented with
\optsshort{a stack}{a sorted linked list}{a binary heap}{an unsorted array}
\end{question}

\begin{question}{M}
Given below are two statements.\\
Statement I: Deleting the maximum from an unsorted array priority queue takes O(1) time.\\
Statement II: Inserting into a sorted linked list priority queue takes O(n) time.
\optslong{Both Statement I and Statement II are true}{Both Statement I and Statement II are false}{Statement I is true but Statement II is false}{Statement I is false but Statement II is true}
\end{question}

\begin{question}{M}
A header node in a linked list is
\optslong{a node that stores the list length only}{the last node}{a dummy node at the beginning that simplifies insertion and deletion at the front}{the first data node}
\end{question}

\begin{question}{M}
Polynomials are commonly represented with linked lists in which each node stores
\optslong{a coefficient, an exponent and a pointer to the next term}{two exponents}{the value of the polynomial at a point}{only the coefficient}
\end{question}

\begin{question}{M}
In the Josephus problem with n = 7 people in a circle and every 3rd person eliminated, starting the count from person 1, the survivor is
\optsshort{4}{7}{3}{5}
\end{question}

\begin{question}{L}
Which pointer of the last node of a singly linked list is NULL?
\optsshort{head}{prev}{next}{data}
\end{question}

\begin{question}{M}
The number of pointer fields that must be changed to insert a node between two nodes of a doubly linked list is
\optsshort{2}{6}{3}{4}
\end{question}

\begin{question}{M}
Given below are two statements.\\
Statement I: Arrays have better cache locality than linked lists.\\
Statement II: Linked lists support O(1) access to the ith element.
\optslong{Both Statement I and Statement II are true}{Both Statement I and Statement II are false}{Statement I is true but Statement II is false}{Statement I is false but Statement II is true}
\end{question}

\begin{answers}
\ans{1}{3}{Merge sort suits linked lists because it needs only sequential access.}
\ans{2}{2}{The neighbours are linked to each other, bypassing p; deletion takes O(1).}
\ans{3}{2}{A linked list has no random access, so binary search is inefficient; the middle is found by walking with slow and fast pointers.}
\ans{4}{3}{The fast pointer moves two steps and the slow pointer one; they meet if and only if there is a cycle.}
\ans{5}{1}{Three pointers (previous, current, next) are moved along the list.}
\ans{6}{1}{Deleting the last node needs its predecessor, which requires a traversal.}
\ans{7}{4}{The recursive call is made before printing, so the values are printed in reverse order.}
\ans{8}{1}{It gives a doubly linked list with one pointer field per node.}
\ans{9}{3}{Nodes are allocated individually and need not be contiguous, which is exactly why lists are flexible.}
\ans{10}{1}{With 0-based indexing they are (i $-$ 1)/2, 2i + 1 and 2i + 2.}
\ans{11}{3}{This is the CLRS example of BUILD-MAX-HEAP.}
\ans{12}{4}{1 is moved to the root and swapped with 14, then 8, then 4.}
\ans{13}{3}{After inserting 30: [30, 20, 15, 10]; 40 is added at index 5 and sifts up past 20 and 30.}
\ans{14}{1}{Most nodes are near the bottom and need little sifting; the total work is bounded by a convergent series.}
\ans{15}{4}{Bottom-up heap construction is linear; heap sort performs n delete-max operations.}
\ans{16}{2}{The minimum must be a leaf, and a heap has $\lceil$n/2$\rceil$ leaves.}
\ans{17}{3}{Nodes $\lfloor$n/2$\rfloor$ + 1 to n are leaves.}
\ans{18}{3}{5 is the root; the left subtree has 3 nodes and the right 1: C(4, 3) $\times$ 2 $\times$ 1 = 8.}
\ans{19}{3}{C(6, 3) $\times$ 2 $\times$ 2 = 80 (each 3-node subheap can be arranged in 2 ways).}
\ans{20}{4}{A heap is a complete binary tree.}
\ans{21}{4}{Heap sort is not stable, because the swaps move equal keys past each other.}
\ans{22}{1}{Each of the n extractions and insertions costs O(log k).}
\ans{23}{2}{Explore only candidates adjacent to removed nodes, keeping at most O(k) items in the auxiliary heap.}
\ans{24}{1}{Mergeable heap variants.}
\ans{25}{1}{Completeness leaves no gaps in the array.}
\ans{26}{4}{In option 2, 5 > 4 (its child); in option 3, 2 > 1; in option 4, 2 > 0.}
\ans{27}{2}{In option 2, 12 < 13 (its child); in option 3, 12 < 16; in option 4, 14 < 16.}
\ans{28}{3}{A heap gives O(log n) insertion and deletion of the maximum.}
\ans{29}{4}{Deleting the maximum from an unsorted array requires a linear scan to find it.}
\ans{30}{3}{It removes special cases for an empty list.}
\ans{31}{1}{Only non-zero terms are stored.}
\ans{32}{1}{Order of elimination: 3, 6, 2, 7, 5, 1; person 4 survives.}
\ans{33}{3}{In a circular list it points back to the first node instead.}
\ans{34}{4}{The new node's prev and next, the predecessor's next, and the successor's prev.}
\ans{35}{3}{Linked list access by position takes O(i).}
\end{answers}

\section{Trees, Binary Trees, Traversals, Threaded Binary Trees and Binary Search Trees}

\begin{question}{H}
A binary tree has preorder traversal A B D E G C F H and inorder traversal D B G E A C H F. Its postorder traversal is
\optsshort{D E G B H F C A}{G E D B H F C A}{D G E B H F C A}{D G E B F H C A}
\end{question}

\begin{question}{M}
A binary tree has preorder A B D H E I C F G and inorder H D B I E A F C G. Its postorder traversal is
\optsshort{H D I E B F G C A}{H I D E B F G C A}{H D I E B G F C A}{D H I E B F G C A}
\end{question}

\begin{question}{M}
A binary tree has preorder F B A D C E G I H and inorder A B C D E F G H I. Its level-order traversal is
\optsshort{F B A D C E G I H}{F B G A I D C E H}{F G B I A D H C E}{F B G A D I C E H}
\end{question}

\begin{question}{H}
Given below are two statements.\\
Statement I: The preorder and postorder traversals together always uniquely determine a binary tree.\\
Statement II: The preorder and inorder traversals together uniquely determine a binary tree with distinct keys.
\optslong{Both Statement I and Statement II are true}{Both Statement I and Statement II are false}{Statement I is true but Statement II is false}{Statement I is false but Statement II is true}
\end{question}

\begin{question}{H}
A binary tree has 20 leaves. The number of nodes with exactly two children is
\optsshort{20}{21}{19}{10}
\end{question}

\begin{question}{H}
The maximum number of nodes in a binary tree of height h (the root at height 0) is
\optsshort{2\textsuperscript{h}}{2\textsuperscript{h+1} $-$ 1}{2\textsuperscript{h} $-$ 1}{2\textsuperscript{h+1}}
\end{question}

\begin{question}{H}
The number of structurally different binary trees with 4 nodes is
\optsshort{16}{10}{14}{12}
\end{question}

\begin{question}{H}
The number of different binary search trees that can be built from the keys 1, 2, 3, 4, 5 is
\optsshort{120}{60}{42}{24}
\end{question}

\begin{question}{M}
The number of different binary trees that can be formed with 3 distinct labelled nodes is
\optsshort{6}{15}{5}{30}
\end{question}

\begin{question}{H}
A full (strictly) binary tree has 15 internal nodes. Its total number of nodes is
\optsshort{30}{31}{32}{29}
\end{question}

\begin{question}{H}
In a complete ternary tree (every internal node has exactly 3 children) with 10 internal nodes, the number of leaves is
\optsshort{31}{21}{30}{20}
\end{question}

\begin{question}{M}
A binary tree with n nodes has how many NULL child pointers when represented with left and right pointers?
\optsshort{n $-$ 1}{n + 1}{2n}{n}
\end{question}

\begin{question}{H}
In a threaded binary tree, the NULL left and right pointers are replaced by pointers to the
\optslong{inorder predecessor and inorder successor}{leftmost and rightmost leaves}{parent and root}{preorder predecessor and successor}
\end{question}

\begin{question}{H}
The keys 50, 15, 62, 5, 20, 58, 91, 3, 8, 37, 60, 24 are inserted in that order into an empty BST. Its postorder traversal is
\optslong{3 8 5 20 24 37 15 58 60 91 62 50}{3 5 8 15 20 24 37 50 58 60 62 91}{3 8 5 24 37 20 15 60 58 91 62 50}{50 15 5 3 8 20 37 24 62 58 60 91}
\end{question}

\begin{question}{M}
The keys 50, 15, 62, 5, 20, 58, 91, 3, 8, 37, 60, 24 are inserted in that order into an empty BST. The height of the tree (the root at height 0) is
\optsshort{3}{4}{6}{5}
\end{question}

\begin{question}{H}
The preorder traversal of a BST is 30 20 10 15 25 23 39 35 42. Its postorder traversal is
\optsshort{15 10 23 25 20 35 42 39 30}{15 23 10 25 20 35 42 39 30}{15 10 25 23 20 42 35 39 30}{10 15 20 23 25 30 35 39 42}
\end{question}

\begin{question}{H}
The inorder traversal of a binary search tree gives
\optslong{the keys in reverse insertion order}{the keys in sorted (non-decreasing) order}{the keys in random order}{the keys in level order}
\end{question}

\begin{question}{H}
Given below are two statements.\\
Statement I: Searching a BST with n nodes takes O(log n) time in the worst case.\\
Statement II: Inserting the keys in sorted order into an empty BST produces a tree of height n $-$ 1.
\optslong{Both Statement I and Statement II are true}{Both Statement I and Statement II are false}{Statement I is true but Statement II is false}{Statement I is false but Statement II is true}
\end{question}

\begin{question}{H}
When deleting a node with two children from a BST, the node is replaced by its
\optslong{left child always}{root}{inorder successor (or inorder predecessor)}{parent}
\end{question}

\begin{question}{M}
The keys 7, 5, 1, 8, 3, 6, 0, 9, 4, 2 are inserted in that order into an empty BST. The number of leaf nodes is
\optsshort{5}{7}{6}{4}
\end{question}

\begin{question}{H}
Match List I with List II.\\
List I (Traversal): A. Preorder  B. Inorder  C. Postorder  D. Level order\\
List II (Order of visiting): I. Left, right, root  II. Root, left, right  III. Level by level, left to right, using a queue  IV. Left, root, right
\optsshort{A-IV, B-II, C-I, D-III}{A-II, B-IV, C-III, D-I}{A-II, B-IV, C-I, D-III}{A-II, B-I, C-IV, D-III}
\end{question}

\begin{question}{M}
For an expression tree of (a + b) * (c $-$ d), the postorder traversal gives
\optsshort{a b + c d $-$ *}{a + b * c $-$ d}{+ a b * $-$ c d}{* + a b $-$ c d}
\end{question}

\begin{question}{M}
Given below are two statements: one is labelled as Assertion A and the other is labelled as Reason R.\\
Assertion A: Recursive traversal of a binary tree of height h needs only O(1) extra space.\\
Reason R: The recursion stack holds one frame per level along the current path.
\optslong{Both A and R are true and R is the correct explanation of A}{Both A and R are true but R is NOT the correct explanation of A}{A is true but R is false}{A is false but R is true}
\end{question}

\begin{question}{M}
The minimum height of a binary tree with n nodes (the root at height 0) is
\optsshort{n $-$ 1}{log2 n squared}{$\lceil$log2(n + 1)$\rceil$ $-$ 1}{$\lfloor$n/2$\rfloor$}
\end{question}

\begin{question}{H}
Match List I with List II.\\
List I (Binary tree type): A. Full binary tree  B. Complete binary tree  C. Perfect binary tree  D. Skewed binary tree\\
List II (Property): I. All levels filled except possibly the last, which is filled from the left  II. Every node has 0 or 2 children  III. Every internal node has two children and all leaves are at the same level  IV. Every node has only one child, except the single leaf
\optsshort{A-II, B-I, C-III, D-IV}{A-II, B-III, C-I, D-IV}{A-II, B-I, C-IV, D-III}{A-I, B-II, C-III, D-IV}
\end{question}

\begin{question}{M}
A general tree is converted to a binary tree using the left-child, right-sibling representation. In the resulting binary tree, the root
\optsshort{has no right child}{has no left child}{is a leaf}{has two children}
\end{question}

\begin{question}{M}
A forest of trees can be represented as a single binary tree by
\optslong{it cannot be done}{joining the trees through their leaves}{deleting all leaves}{connecting the roots of the trees as siblings in the left-child, right-sibling representation}
\end{question}

\begin{question}{M}
The number of edges in a tree with n nodes is
\optsshort{n $-$ 1}{n}{n + 1}{2n}
\end{question}

\begin{question}{M}
A tree has 4 nodes of degree 3 and 3 nodes of degree 2; all other nodes are leaves (degree 1). The number of leaves is
\optsshort{5}{7}{6}{4}
\end{question}

\begin{question}{M}
Given below are two statements.\\
Statement I: Every binary search tree is a heap.\\
Statement II: A heap stores the keys so that an inorder traversal yields them in sorted order.
\optslong{Both Statement I and Statement II are true}{Both Statement I and Statement II are false}{Statement I is true but Statement II is false}{Statement I is false but Statement II is true}
\end{question}

\begin{question}{M}
Which traversal of a binary tree can be used to delete (free) all its nodes safely?
\optsshort{Level order}{Preorder}{Inorder}{Postorder}
\end{question}

\begin{question}{M}
The time complexity of finding the height of a binary tree with n nodes is
\optsshort{O(n)}{O(n\textsuperscript{2})}{O(log n)}{O(1)}
\end{question}

\begin{question}{M}
What does the following function return for a binary tree?\\
int f(Node *t) \{ if (t == NULL) return 0; if (t->left == NULL \&\& t->right == NULL) return 1; return f(t->left) + f(t->right); \}
\optsshort{The number of internal nodes}{The number of nodes}{The height}{The number of leaves}
\end{question}

\begin{question}{M}
In a binary search tree, the minimum key is found by
\optslong{following right pointers from the root}{following left pointers from the root until a node with no left child}{checking all leaves}{a level-order traversal}
\end{question}

\begin{question}{M}
The inorder successor of a node with a right child in a BST is
\optslong{its parent}{the root}{the leftmost node of its right subtree}{its right child always}
\end{question}

\begin{question}{M}
Given below are two statements: one is labelled as Assertion A and the other is labelled as Reason R.\\
Assertion A: A threaded binary tree can be traversed in inorder without a stack.\\
Reason R: Threads point to the parent of each node.
\optslong{Both A and R are true and R is the correct explanation of A}{Both A and R are true but R is NOT the correct explanation of A}{A is true but R is false}{A is false but R is true}
\end{question}

\begin{question}{M}
An expression tree for an arithmetic expression has operators at
\optsshort{the internal nodes}{no nodes}{only the root}{the leaves}
\end{question}

\begin{question}{M}
In a complete binary tree with n nodes stored in an array (1-indexed), node i is a leaf if
\optsshort{i = 1}{i > n/2}{i < n/2}{i is even}
\end{question}

\begin{question}{L}
The maximum number of nodes at level k of a binary tree (the root at level 0) is
\optsshort{2k}{k\textsuperscript{2}}{k}{2\textsuperscript{k}}
\end{question}

\begin{question}{M}
Given below are two statements.\\
Statement I: The level-order traversal of a binary tree uses a stack.\\
Statement II: The postorder traversal of a binary tree visits the root last.
\optslong{Both Statement I and Statement II are true}{Both Statement I and Statement II are false}{Statement I is true but Statement II is false}{Statement I is false but Statement II is true}
\end{question}

\begin{answers}
\ans{1}{3}{A is the root; the left subtree (D B G E) has root B; the right subtree (C H F) has root C with right child F whose left child is H.}
\ans{2}{1}{A is the root; B has left child D (with left child H) and right child E (with left child I); C has children F and G.}
\ans{3}{4}{F is the root with children B and G; B has A and D; D has C and E; G has right child I, which has left child H.}
\ans{4}{4}{Preorder and postorder cannot distinguish a single left child from a single right child, unless the tree is full.}
\ans{5}{3}{In any binary tree, n0 = n2 + 1, so n2 = 19.}
\ans{6}{2}{Level i has at most 2\textsuperscript{i} nodes: 1 + 2 + ... + 2\textsuperscript{h} = 2\textsuperscript{h+1} $-$ 1.}
\ans{7}{3}{It is the Catalan number C(4) = 14.}
\ans{8}{3}{It is the Catalan number C(5) = 42.}
\ans{9}{4}{n! $\times$ C(n) = 6 $\times$ 5 = 30 (5 shapes, each labelled in 3! ways).}
\ans{10}{2}{A full binary tree with i internal nodes has i + 1 leaves: 15 + 16 = 31.}
\ans{11}{2}{For a full k-ary tree, leaves = (k $-$ 1) $\times$ i + 1 = 2 $\times$ 10 + 1 = 21.}
\ans{12}{2}{2n pointer fields, of which n $-$ 1 point to children.}
\ans{13}{1}{Threads allow inorder traversal without a stack or recursion.}
\ans{14}{3}{Option 2 is the inorder and option 3 the preorder of the same tree.}
\ans{15}{2}{The longest path is 50 $\rightarrow$ 15 $\rightarrow$ 20 $\rightarrow$ 37 $\rightarrow$ 24.}
\ans{16}{1}{Rebuild the BST from the preorder sequence: 30 has left subtree 20 (10 $\rightarrow$ 15, 25 $\rightarrow$ 23) and right subtree 39 (35, 42).}
\ans{17}{2}{This is why the inorder traversal of a BST can be used for sorting.}
\ans{18}{4}{The worst case for search is O(n), which occurs exactly when the tree degenerates into a chain, as in Statement II.}
\ans{19}{3}{The successor is the minimum of the right subtree and has at most one child.}
\ans{20}{1}{The leaves are 0, 2, 4, 6 and 9.}
\ans{21}{3}{Depth-first traversals use a stack (or recursion); level order uses a queue.}
\ans{22}{1}{Postorder gives the postfix form; preorder gives the prefix form.}
\ans{23}{4}{Because of R, recursive traversal needs O(h) space; Morris traversal achieves O(1).}
\ans{24}{3}{A tree of height h holds at most 2\textsuperscript{h+1} $-$ 1 nodes.}
\ans{25}{1}{Definitions vary slightly between textbooks; these are the common ones.}
\ans{26}{1}{The root of a tree has no siblings.}
\ans{27}{4}{The roots of successive trees become right siblings.}
\ans{28}{1}{A tree is connected and acyclic.}
\ans{29}{3}{Sum of degrees = 2(n $-$ 1): 12 + 6 + L = 2(7 + L $-$ 1), giving L = 6.}
\ans{30}{2}{A BST orders keys left-to-right; a heap orders them only along root-to-leaf paths.}
\ans{31}{4}{Children are freed before their parent.}
\ans{32}{1}{Every node is visited once.}
\ans{33}{4}{It returns 1 for each node with no children and adds the counts of the subtrees.}
\ans{34}{2}{The maximum is found by following right pointers.}
\ans{35}{3}{If there is no right child, the successor is the lowest ancestor whose left subtree contains the node.}
\ans{36}{3}{Threads point to the inorder predecessor and successor, not to the parent.}
\ans{37}{1}{Operands are at the leaves.}
\ans{38}{2}{Its children 2i and 2i + 1 would be beyond n.}
\ans{39}{4}{Each level can at most double the number of nodes.}
\ans{40}{4}{Level-order traversal uses a queue.}
\end{answers}

\section{Balanced Search Trees (AVL, red-black, B, B+, B* trees) and Data Structures for Sets}

\begin{question}{H}
The keys 1, 2, 3, 4, 5, 6, 7 are inserted in that order into an empty AVL tree. The root of the final tree is
\optsshort{3}{7}{4}{1}
\end{question}

\begin{question}{H}
The keys 10, 20, 30, 40, 50, 25 are inserted in that order into an empty AVL tree. The preorder traversal of the result is
\optsshort{40 20 10 30 25 50}{20 10 30 25 40 50}{30 20 10 25 40 50}{30 20 25 10 40 50}
\end{question}

\begin{question}{M}
The keys 50, 20, 60, 10, 8, 15, 32, 46, 11, 48 are inserted in that order into an empty AVL tree. The root of the final tree is
\optsshort{10}{20}{50}{46}
\end{question}

\begin{question}{H}
The minimum number of nodes in an AVL tree of height 4 (a single node has height 0) is
\optsshort{15}{12}{20}{7}
\end{question}

\begin{question}{H}
The minimum number of nodes in an AVL tree of height 5 (a single node has height 0) is
\optsshort{16}{33}{20}{12}
\end{question}

\begin{question}{H}
The balance factor of every node in an AVL tree must be
\optsshort{$-$2 to +2}{positive}{0}{$-$1, 0 or +1}
\end{question}

\begin{question}{H}
Match List I with List II.\\
List I (Imbalance case at node z after insertion): A. Left-Left  B. Right-Right  C. Left-Right  D. Right-Left\\
List II (Fix): I. Single left rotation at z  II. Single right rotation at z  III. Right rotation at z's right child, then left rotation at z  IV. Left rotation at z's left child, then right rotation at z
\optsshort{A-II, B-I, C-III, D-IV}{A-I, B-II, C-IV, D-III}{A-II, B-I, C-IV, D-III}{A-II, B-IV, C-I, D-III}
\end{question}

\begin{question}{H}
Given below are two statements.\\
Statement I: The height of an AVL tree with n nodes is O(log n).\\
Statement II: Insertion into an AVL tree may require O(log n) rotations.
\optslong{Both Statement I and Statement II are true}{Both Statement I and Statement II are false}{Statement I is true but Statement II is false}{Statement I is false but Statement II is true}
\end{question}

\begin{question}{M}
The maximum height of an AVL tree with n nodes is approximately
\optsshort{1.44 log2 n}{2 log2 n}{log2 n}{n/2}
\end{question}

\begin{question}{H}
Match List I with List II.\\
List I (Red-black tree property): A. Root property  B. Red property  C. Black-height property  D. Leaf property\\
List II (Statement): I. Every NIL leaf is black  II. Every path from a node to its descendant NIL leaves has the same number of black nodes  III. The root is black  IV. The children of a red node are black
\optsshort{A-IV, B-III, C-II, D-I}{A-III, B-IV, C-I, D-II}{A-III, B-IV, C-II, D-I}{A-III, B-II, C-IV, D-I}
\end{question}

\begin{question}{H}
The height of a red-black tree with n internal nodes is at most
\optsshort{1.44 log2 n}{2 log2(n + 1)}{log2 n}{n}
\end{question}

\begin{question}{M}
Given below are two statements: one is labelled as Assertion A and the other is labelled as Reason R.\\
Assertion A: Red-black trees usually need fewer rotations on insertion and deletion than AVL trees.\\
Reason R: Red-black trees are more strictly balanced than AVL trees.
\optslong{Both A and R are true and R is the correct explanation of A}{Both A and R are true but R is NOT the correct explanation of A}{A is true but R is false}{A is false but R is true}
\end{question}

\begin{question}{H}
In a B-tree of order m (maximum m children per node), every internal node other than the root has at least
\optsshort{m/4 children}{m $-$ 1 children}{$\lceil$m/2$\rceil$ children}{2 children}
\end{question}

\begin{question}{H}
In a B-tree of order 5, the maximum and minimum numbers of keys in a non-root node are
\optsshort{4 and 1}{5 and 3}{4 and 2}{5 and 2}
\end{question}

\begin{question}{H}
Given below are two statements.\\
Statement I: In a B-tree, all leaves are at the same level.\\
Statement II: A B-tree grows in height when the root splits.
\optslong{Both Statement I and Statement II are true}{Both Statement I and Statement II are false}{Statement I is true but Statement II is false}{Statement I is false but Statement II is true}
\end{question}

\begin{question}{H}
Match List I with List II.\\
List I (Tree): A. B tree  B. B+ tree  C. B* tree  D. 2-3 tree\\
List II (Feature): I. Nodes are kept at least two-thirds full by redistributing with siblings before splitting  II. Keys and data may be in any node  III. A B-tree of order 3  IV. All data is in the leaves, which are linked for sequential access
\optsshort{A-II, B-I, C-IV, D-III}{A-IV, B-II, C-I, D-III}{A-II, B-IV, C-III, D-I}{A-II, B-IV, C-I, D-III}
\end{question}

\begin{question}{M}
The keys 1, 2, 3, 4, 5 are inserted in that order into an initially empty B-tree of order 3 (2-3 tree). The root of the final tree contains
\optsshort{2 and 4}{3}{4}{2}
\end{question}

\begin{question}{M}
A splay tree performs all of the following operations in O(log n) amortized time. What does a splay operation do?
\optslong{Moves the accessed node to the root using rotations}{Rebalances the tree using colours}{Deletes the accessed node}{Converts the tree to a heap}
\end{question}

\begin{question}{H}
In the disjoint-set (union-find) data structure with union by rank and path compression, the amortized time per operation is
\optslong{O(log log n)}{O(1)}{O(log n)}{O($\alpha$(n)), where $\alpha$ is the inverse Ackermann function}
\end{question}

\begin{question}{H}
Match List I with List II.\\
List I (Union-find heuristic): A. Union by rank  B. Path compression  C. Linked-list representation with weighted union\\
List II (Effect): I. A sequence of m operations on n elements takes O(m + n log n)  II. Makes every node on the find path point directly to the root  III. Attaches the root of the shorter tree under the root of the taller tree
\optsshort{A-II, B-III, C-I}{A-III, B-I, C-II}{A-III, B-II, C-I}{A-I, B-II, C-III}
\end{question}

\begin{question}{M}
Union-find is used in Kruskal's algorithm to
\optslong{check whether two vertices are already in the same component}{sort the edges}{find shortest paths}{store the adjacency list}
\end{question}

\begin{question}{M}
With union by rank alone (no path compression), the height of any tree in a disjoint-set forest of n elements is at most
\optsshort{log2 n}{n/2}{$\surd$n}{1}
\end{question}

\begin{question}{M}
Given below are two statements.\\
Statement I: Testing membership in a bit-vector set takes O(n) time.\\
Statement II: A set over a small universe \{0, ..., n $-$ 1\} can be represented as a bit vector, with union and intersection done by bitwise OR and AND.
\optslong{Both Statement I and Statement II are true}{Both Statement I and Statement II are false}{Statement I is true but Statement II is false}{Statement I is false but Statement II is true}
\end{question}

\begin{question}{M}
In a B+ tree, range queries are efficient because
\optslong{the tree has height 1}{all internal nodes store data records}{keys are stored in the root}{leaf nodes are linked in key order, so the range is scanned sequentially after one search}
\end{question}

\begin{question}{M}
The keys 3, 2, 1, 4, 5, 6, 7, 16, 15, 14 are inserted in that order into an empty AVL tree. The root of the resulting tree is
\optsshort{4}{3}{15}{7}
\end{question}

\begin{question}{M}
The number of rotations needed in the worst case to delete a node from an AVL tree with n nodes is
\optsshort{1}{O(n)}{O(log n)}{2}
\end{question}

\begin{question}{M}
Given below are two statements: one is labelled as Assertion A and the other is labelled as Reason R.\\
Assertion A: B-trees are preferred for indexes stored on disk.\\
Reason R: B-trees were introduced by Bayer and McCreight.
\optslong{Both A and R are true and R is the correct explanation of A}{Both A and R are true but R is NOT the correct explanation of A}{A is true but R is false}{A is false but R is true}
\end{question}

\begin{question}{M}
Which of the following trees guarantees O(log n) worst-case time for search, insert and delete?\\
A. AVL tree  B. Red-black tree  C. Binary search tree  D. B-tree\\
Choose the correct answer from the options given below:
\optsshort{B and C only}{A and C only}{A, B, C and D}{A, B and D only}
\end{question}

\begin{question}{M}
A trie (prefix tree) is most suitable for
\optslong{implementing priority queues}{sorting integers}{graph traversal}{storing strings for fast prefix search and dictionary look-up}
\end{question}

\begin{question}{M}
A segment tree over an array of n elements supports range-sum queries and point updates each in
\optsshort{O(n log n)}{O(n)}{O(log n)}{O(1)}
\end{question}

\begin{question}{L}
AVL trees are named after
\optsshort{Adelson-Velsky and Landis}{Aho, Vazquez and Lam}{Adams, Venn and Lee}{Allen, Valiant and Lin}
\end{question}

\begin{question}{M}
Given below are two statements.\\
Statement I: Every red-black tree is an AVL tree.\\
Statement II: The height of a red-black tree with n nodes can be $\Theta$(n).
\optslong{Both Statement I and Statement II are true}{Both Statement I and Statement II are false}{Statement I is true but Statement II is false}{Statement I is false but Statement II is true}
\end{question}

\begin{question}{M}
In a B-tree of order m with n keys, the height is
\optsshort{O(m log n)}{O(n)}{O(log\_m n)}{O(n/m)}
\end{question}

\begin{answers}
\ans{1}{3}{Rotations keep the tree balanced, and the result is the perfect tree with root 4.}
\ans{2}{3}{Inserting 25 unbalances 20 (or the root), and a double (right-left) rotation makes 30 the root.}
\ans{3}{2}{After several rotations, 20 becomes the root, with left subtree 10 and right subtree 50.}
\ans{4}{2}{N(h) = N(h $-$ 1) + N(h $-$ 2) + 1 with N(0) = 1 and N(1) = 2: N(2) = 4, N(3) = 7, N(4) = 12.}
\ans{5}{3}{N(5) = N(4) + N(3) + 1 = 12 + 7 + 1 = 20.}
\ans{6}{4}{Balance factor = height of the left subtree $-$ height of the right subtree.}
\ans{7}{3}{Single rotations fix outside cases; double rotations fix inside cases.}
\ans{8}{3}{Insertion needs at most one single or double rotation; deletion may need O(log n) rotations.}
\ans{9}{1}{The minimum-node AVL trees are Fibonacci trees, giving about 1.44 log2(n + 2).}
\ans{10}{3}{These properties keep the tree approximately balanced.}
\ans{11}{2}{The longest path is at most twice the shortest because red nodes cannot be consecutive.}
\ans{12}{3}{Red-black trees are less strictly balanced, which is why they rebalance less often.}
\ans{13}{3}{It therefore has at least $\lceil$m/2$\rceil$ $-$ 1 keys; the root needs at least 2 children unless it is a leaf.}
\ans{14}{3}{At most m $-$ 1 = 4 keys, and at least $\lceil$5/2$\rceil$ $-$ 1 = 2 keys.}
\ans{15}{1}{B-trees grow upward from the root, which keeps them perfectly height-balanced.}
\ans{16}{4}{Variants of the B-tree.}
\ans{17}{1}{Inserting 3 splits [1, 2, 3] with 2 going up; inserting 5 splits [3, 4, 5] with 4 going up.}
\ans{18}{1}{Frequently accessed items stay near the root.}
\ans{19}{4}{$\alpha$(n) $\leq$ 4 for all practical values of n.}
\ans{20}{3}{Together, union by rank and path compression give the inverse-Ackermann bound.}
\ans{21}{1}{An edge is added only if its endpoints are in different sets.}
\ans{22}{1}{A tree of rank r has at least 2\textsuperscript{r} nodes.}
\ans{23}{4}{Membership is a single bit test, O(1).}
\ans{24}{4}{This is a key advantage over B-trees for database indexes.}
\ans{25}{1}{After inserting 14, a double rotation at 6 makes 7 the right child of root 4.}
\ans{26}{3}{Rebalancing may propagate up to the root.}
\ans{27}{2}{Both are true, but B-trees suit disks because a node with many keys fills a disk block and keeps the tree shallow.}
\ans{28}{4}{An unbalanced BST can degenerate into a list.}
\ans{29}{4}{Search time depends on the key length, not the number of keys.}
\ans{30}{3}{It needs O(n) space.}
\ans{31}{1}{They were proposed in 1962 as the first self-balancing BSTs.}
\ans{32}{2}{Red-black trees are less strictly balanced than AVL trees, but their height is still at most 2 log2(n + 1).}
\ans{33}{3}{Each level multiplies the number of keys by at least about m/2.}
\end{answers}

\section{Hashing (hash functions, collision resolution, load factor, analysis)}

\begin{question}{H}
The keys 43, 165, 62, 123, 142 are inserted in that order into a hash table of size 11 with h(k) = k mod 11 and linear probing. The key 142 is stored at index
\optsshort{10}{3}{2}{1}
\end{question}

\begin{question}{H}
The keys 89, 18, 49, 58, 69 are inserted in that order into a table of size 10 with h(k) = k mod 10 and linear probing. The keys 58 and 69 are stored at indices
\optsshort{1 and 2}{2 and 3}{0 and 1}{8 and 9}
\end{question}

\begin{question}{H}
The keys 89, 18, 49, 58, 69 are inserted in that order into a table of size 10 with h(k) = k mod 10 and quadratic probing (h(k) + i\textsuperscript{2}) mod 10. The keys 58 and 69 are stored at indices
\optsshort{2 and 4}{1 and 2}{2 and 3}{3 and 0}
\end{question}

\begin{question}{H}
The keys 89, 18, 49, 58, 69 are inserted in that order into a table of size 10 with h1(k) = k mod 10, h2(k) = 7 $-$ (k mod 7) and double hashing (h1 + i$\cdot$h2) mod 10. The key 49 is stored at index
\optsshort{6}{9}{0}{3}
\end{question}

\begin{question}{M}
The keys 12, 18, 13, 2, 3, 23, 5, 15 are inserted into a table of size 10 with h(k) = k mod 10 and linear probing. The key 15 ends up at index
\optsshort{5}{8}{7}{9}
\end{question}

\begin{question}{H}
The keys 15, 11, 27, 8, 12, 32, 22, 37 are inserted into a table of size 7 using h(k) = k mod 7 and separate chaining. The length of the longest chain is
\optsshort{5}{3}{4}{2}
\end{question}

\begin{question}{H}
Match List I with List II.\\
List I (Collision resolution): A. Linear probing  B. Quadratic probing  C. Double hashing  D. Separate chaining\\
List II (Characteristic): I. Avoids both primary and secondary clustering  II. Suffers from primary clustering  III. Colliding keys are stored in a linked list at the slot  IV. Avoids primary clustering but suffers from secondary clustering
\optsshort{A-II, B-I, C-IV, D-III}{A-IV, B-II, C-I, D-III}{A-II, B-IV, C-I, D-III}{A-II, B-IV, C-III, D-I}
\end{question}

\begin{question}{H}
In open addressing with uniform hashing and load factor $\alpha$ < 1, the expected number of probes in an unsuccessful search is at most
\optsshort{$\alpha$}{1/(1 $-$ $\alpha$)}{1 + $\alpha$}{log(1/$\alpha$)}
\end{question}

\begin{question}{M}
With open addressing and uniform hashing, the expected number of probes in an unsuccessful search when the table is 90\% full is at most
\optsshort{9}{10}{5}{2}
\end{question}

\begin{question}{M}
With open addressing and uniform hashing, the expected number of probes in a successful search at load factor $\alpha$ = 0.5 is about
\optsshort{1.5}{2}{2.5}{1.39}
\end{question}

\begin{question}{H}
In separate chaining with load factor $\alpha$ = n/m and simple uniform hashing, the expected time for an unsuccessful search is
\optsshort{$\Theta$(n)}{$\Theta$(1)}{$\Theta$($\alpha$\textsuperscript{2})}{$\Theta$(1 + $\alpha$)}
\end{question}

\begin{question}{H}
Given below are two statements.\\
Statement I: In open addressing, the load factor can exceed 1.\\
Statement II: In separate chaining, the load factor can exceed 1.
\optslong{Both Statement I and Statement II are true}{Both Statement I and Statement II are false}{Statement I is true but Statement II is false}{Statement I is false but Statement II is true}
\end{question}

\begin{question}{H}
Given below are two statements: one is labelled as Assertion A and the other is labelled as Reason R.\\
Assertion A: In open addressing, a deleted slot is usually marked as ``deleted'' instead of being made empty.\\
Reason R: Emptying the slot could break the probe sequence and make later keys unreachable by search.
\optslong{Both A and R are true and R is the correct explanation of A}{Both A and R are true but R is NOT the correct explanation of A}{A is true but R is false}{A is false but R is true}
\end{question}

\begin{question}{H}
In the division method h(k) = k mod m, a good choice of m is
\optslong{any even number}{a prime not too close to a power of 2}{a power of 10}{a power of 2}
\end{question}

\begin{question}{M}
The multiplication method of hashing computes h(k) = $\lfloor$m (kA mod 1)$\rfloor$ with 0 < A < 1. Knuth suggests A $\approx$
\optslong{0.5}{0.75}{0.618 (the golden-ratio conjugate)}{0.333}
\end{question}

\begin{question}{M}
Three keys are hashed independently and uniformly into a table of 10 slots. The probability that there is at least one collision is
\optsshort{0.28}{0.30}{0.10}{0.72}
\end{question}

\begin{question}{M}
In universal hashing, for any two distinct keys, the probability that they collide under a hash function chosen at random from the family is at most
\optsshort{1/m}{1/n}{1/m\textsuperscript{2}}{1/2}
\end{question}

\begin{question}{M}
Given below are two statements.\\
Statement I: Perfect hashing guarantees O(1) worst-case search time for a static set of keys.\\
Statement II: Perfect hashing typically uses two levels of hashing with universal hash functions.
\optslong{Both Statement I and Statement II are true}{Both Statement I and Statement II are false}{Statement I is true but Statement II is false}{Statement I is false but Statement II is true}
\end{question}

\begin{question}{M}
Rehashing is performed when
\optslong{a search fails}{the hash function returns zero}{a key is deleted}{the load factor exceeds a threshold, so a larger table is allocated and all keys are reinserted}
\end{question}

\begin{question}{M}
Linear probing suffers from primary clustering, which means that
\optslong{keys are stored in linked lists}{long runs of occupied slots form and grow, increasing search times}{keys that hash to the same slot follow the same probe sequence}{the table can never be full}
\end{question}

\begin{question}{M}
Quadratic probing with a prime table size is guaranteed to find an empty slot for insertion if the table is
\optsshort{of any load factor}{at least 90\% full}{completely full}{at most half full}
\end{question}

\begin{question}{M}
Cuckoo hashing guarantees
\optslong{no rehashing ever}{O(log n) lookup}{sorted order of keys}{O(1) worst-case lookup time, since a key is in one of two possible locations}
\end{question}

\begin{question}{M}
Which of the following is NOT a desirable property of a hash function?
\optslong{It uses all of the key}{It is fast to compute}{It distributes keys uniformly over the table}{It maps similar keys to the same slot}
\end{question}

\begin{question}{M}
Given below are two statements: one is labelled as Assertion A and the other is labelled as Reason R.\\
Assertion A: Hash tables are not suitable for range queries.\\
Reason R: A hash table keeps its keys in sorted order.
\optslong{Both A and R are true and R is the correct explanation of A}{Both A and R are true but R is NOT the correct explanation of A}{A is true but R is false}{A is false but R is true}
\end{question}

\begin{question}{M}
A hash table of size 10 uses h(k) = k mod 10 and linear probing. After inserting 21, 31, 41 in that order, how many probes are needed to insert 51?
\optsshort{3}{1}{2}{4}
\end{question}

\begin{question}{M}
In separate chaining, inserting a new key at the head of the chain takes
\optsshort{O(1) time}{O($\alpha$) time}{O(n) time}{O(log n) time}
\end{question}

\begin{question}{M}
In double hashing, the second hash function h2(k) must never evaluate to
\optsshort{an odd number}{a prime}{0}{1}
\end{question}

\begin{question}{M}
Given below are two statements.\\
Statement I: Extendible hashing uses a directory of pointers to buckets.\\
Statement II: Linear hashing requires a directory that doubles on every split.
\optslong{Both Statement I and Statement II are true}{Both Statement I and Statement II are false}{Statement I is true but Statement II is false}{Statement I is false but Statement II is true}
\end{question}

\begin{question}{L}
A Bloom filter can report
\optslong{only exact counts}{false negatives but not false positives}{neither false positives nor false negatives}{false positives but not false negatives}
\end{question}

\begin{question}{M}
The keys 25, 36, 47, 58 are inserted into a table of size 11 with h(k) = k mod 11 and linear probing. The number of collisions (occupied slots probed) is
\optsshort{3}{0}{10}{6}
\end{question}

\begin{question}{M}
In the mid-square hashing method
\optslong{the key is multiplied by 0.618}{the key is split into parts that are added}{the key is squared and the middle digits of the result are used as the address}{the key is divided by the table size}
\end{question}

\begin{answers}
\ans{1}{4}{142 mod 11 = 10 is taken by 43; index 0 is taken by 165; so 142 goes to index 1.}
\ans{2}{1}{49 goes to 0 (after 9); 58 tries 8, 9, 0 and goes to 1; 69 tries 9, 0, 1 and goes to 2.}
\ans{3}{3}{58: 8, 9, then (8 + 4) mod 10 = 2. 69: 9, 0, then (9 + 4) mod 10 = 3.}
\ans{4}{1}{49: h1 = 9 is occupied; h2 = 7 $-$ 0 = 7; (9 + 7) mod 10 = 6.}
\ans{5}{4}{Indices 2--8 are filled by 12, 13, 2, 3, 23, 5 and 18; 15 probes 5, 6, 7, 8 and lands at 9.}
\ans{6}{2}{15, 8 and 22 all hash to 1.}
\ans{7}{3}{Standard properties of collision-resolution methods.}
\ans{8}{2}{For $\alpha$ = 0.5 this is 2 probes; for $\alpha$ = 0.9 it is 10.}
\ans{9}{2}{1/(1 $-$ 0.9) = 10.}
\ans{10}{4}{(1/$\alpha$) ln(1/(1 $-$ $\alpha$)) = 2 ln 2 $\approx$ 1.39.}
\ans{11}{4}{The whole chain of length $\alpha$ on average is scanned, plus the cost of hashing.}
\ans{12}{4}{Open addressing stores every key in the table itself, so n $\leq$ m.}
\ans{13}{1}{Such tombstones are skipped by searches but reused by insertions.}
\ans{14}{2}{A power of 2 would use only the low-order bits of the key.}
\ans{15}{3}{($\surd$5 $-$ 1)/2 $\approx$ 0.618 works well; the value of m is not critical in this method.}
\ans{16}{1}{P(no collision) = (10 $\times$ 9 $\times$ 8) / 10\textsuperscript{3} = 0.72, so P(collision) = 0.28.}
\ans{17}{1}{This guarantees good expected performance for any input.}
\ans{18}{1}{The second-level table for a slot with n\_j keys has size n\_j\textsuperscript{2}, which makes collisions avoidable.}
\ans{19}{4}{Doubling the table keeps the amortized insertion cost O(1).}
\ans{20}{2}{Secondary clustering refers to keys with the same home slot sharing a probe sequence.}
\ans{21}{4}{The first $\lceil$m/2$\rceil$ probe positions are distinct when m is prime.}
\ans{22}{4}{Insertion may displace keys and occasionally requires rehashing.}
\ans{23}{4}{Similar keys should be spread out to avoid clustering.}
\ans{24}{3}{A good hash function scatters keys, so neighbouring keys are not stored near each other.}
\ans{25}{4}{51 probes slots 1, 2, 3 (taken by 21, 31, 41) and then 4.}
\ans{26}{1}{Checking for duplicates first would cost O(1 + $\alpha$).}
\ans{27}{3}{A zero step would probe the same slot forever; h2(k) should also be relatively prime to m.}
\ans{28}{3}{Linear hashing needs no directory; buckets are split in a fixed linear order.}
\ans{29}{4}{It is a space-efficient probabilistic set-membership structure.}
\ans{30}{4}{All four hash to 3: 36 makes 1 probe of an occupied slot, 47 makes 2 and 58 makes 3, total 6.}
\ans{31}{3}{Folding splits the key into parts that are added together.}
\end{answers}

\section{Sorting and Searching}

\begin{question}{H}
Match List I with List II.\\
List I (Sorting algorithm): A. Quick sort  B. Merge sort  C. Insertion sort  D. Heap sort\\
List II (Worst-case time): I. O(n log n), but not in place  II. O(n\textsuperscript{2}), but O(n log n) on average  III. O(n log n) and in place  IV. O(n\textsuperscript{2}), but O(n) on already sorted input
\optsshort{A-II, B-IV, C-I, D-III}{A-II, B-I, C-IV, D-III}{A-II, B-I, C-III, D-IV}{A-I, B-II, C-IV, D-III}
\end{question}

\begin{question}{H}
Which of the following sorting algorithms are stable in their standard implementations?\\
A. Insertion sort  B. Merge sort  C. Quick sort  D. Heap sort  E. Counting sort\\
Choose the correct answer from the options given below:
\optsshort{A, B and E only}{B, C and E only}{A, C and D only}{A, B, C and E only}
\end{question}

\begin{question}{H}
The lower bound on the number of comparisons needed by any comparison-based sorting algorithm in the worst case is
\optsshort{$\Omega$(log n)}{$\Omega$(n log n)}{$\Omega$(n\textsuperscript{2})}{$\Omega$(n)}
\end{question}

\begin{question}{H}
The minimum number of comparisons needed in the worst case to sort 5 elements is at least
\optsshort{10}{5}{7}{8}
\end{question}

\begin{question}{H}
Quick sort, using the first or last element as the pivot, takes O(n\textsuperscript{2}) time when the input is
\optsshort{already sorted}{random}{all distinct}{of even length}
\end{question}

\begin{question}{H}
The CLRS (Lomuto) PARTITION procedure, using the last element as the pivot, is applied to [2, 8, 7, 1, 3, 5, 6, 4]. The resulting array is
\optsshort{[1, 2, 3, 4, 5, 6, 7, 8]}{[2, 1, 3, 4, 7, 5, 6, 8]}{[3, 1, 2, 4, 7, 5, 6, 8]}{[2, 1, 3, 4, 8, 7, 5, 6]}
\end{question}

\begin{question}{H}
After the first pass of bubble sort (ascending) on [5, 1, 4, 2, 8], the array is
\optsshort{[1, 5, 4, 2, 8]}{[1, 2, 4, 5, 8]}{[1, 4, 2, 5, 8]}{[5, 1, 2, 4, 8]}
\end{question}

\begin{question}{H}
After two passes of selection sort (ascending) on [64, 25, 12, 22, 11], the array is
\optsshort{[11, 12, 25, 22, 64]}{[11, 25, 12, 22, 64]}{[12, 11, 25, 22, 64]}{[11, 12, 22, 25, 64]}
\end{question}

\begin{question}{H}
The number of element shifts performed by insertion sort on [8, 4, 3, 7, 6, 5, 2, 1] equals the number of inversions, which is
\optsshort{22}{28}{20}{18}
\end{question}

\begin{question}{H}
Given below are two statements.\\
Statement I: Selection sort performs O(n) comparisons in the best case.\\
Statement II: Selection sort performs O(n) swaps in the worst case.
\optslong{Both Statement I and Statement II are true}{Both Statement I and Statement II are false}{Statement I is true but Statement II is false}{Statement I is false but Statement II is true}
\end{question}

\begin{question}{H}
In LSD radix sort of [170, 45, 75, 90, 802, 24, 2, 66], the array after sorting on the units digit is
\optslong{[170, 90, 2, 802, 24, 45, 75, 66]}{[2, 24, 45, 66, 75, 90, 170, 802]}{[802, 2, 24, 45, 66, 170, 75, 90]}{[170, 90, 802, 2, 24, 45, 75, 66]}
\end{question}

\begin{question}{H}
Counting sort sorts n integers in the range 0 to k in
\optsshort{O(nk)}{O(n + k)}{O(k log n)}{O(n log n)}
\end{question}

\begin{question}{M}
Radix sort on n numbers of d digits in base b, using counting sort for each digit, takes
\optsshort{O(d(n + b))}{O(b\textsuperscript{d})}{O(n log n)}{O(n\textsuperscript{d})}
\end{question}

\begin{question}{M}
Bucket sort runs in linear expected time when
\optslong{there is one bucket}{the input is sorted in reverse}{the input is uniformly distributed over a range}{all keys are equal}
\end{question}

\begin{question}{H}
The maximum number of comparisons made by binary search on a sorted array of 1000 elements is
\optsshort{9}{500}{10}{11}
\end{question}

\begin{question}{H}
Given below are two statements: one is labelled as Assertion A and the other is labelled as Reason R.\\
Assertion A: Binary search cannot be applied directly to an unsorted array.\\
Reason R: Binary search discards half of the elements based on a comparison with the middle element, which is valid only if the array is ordered.
\optslong{Both A and R are true and R is the correct explanation of A}{Both A and R are true but R is NOT the correct explanation of A}{A is true but R is false}{A is false but R is true}
\end{question}

\begin{question}{M}
Interpolation search on uniformly distributed sorted data takes, on average,
\optsshort{O(n)}{O(log n)}{O(1)}{O(log log n)}
\end{question}

\begin{question}{H}
Merging two sorted lists of m and n elements requires at most how many comparisons?
\optsshort{m + n}{max(m, n)}{m $\times$ n}{m + n $-$ 1}
\end{question}

\begin{question}{M}
The minimum number of comparisons needed to find both the maximum and minimum of 100 elements is
\optsshort{150}{148}{99}{198}
\end{question}

\begin{question}{M}
The minimum number of comparisons needed to find the second largest of 64 elements is
\optsshort{126}{70}{68}{63}
\end{question}

\begin{question}{H}
The kth smallest element of an unsorted array can be found in worst-case linear time using
\optslong{heap sort}{binary search}{the median-of-medians selection algorithm}{quick sort}
\end{question}

\begin{question}{M}
Match List I with List II.\\
List I (Sorting algorithm): A. Bubble sort  B. Selection sort  C. Merge sort  D. Shell sort\\
List II (Idea): I. Repeatedly picks the minimum of the unsorted part  II. Repeatedly swaps adjacent out-of-order elements  III. Insertion sort on elements a gap apart, with decreasing gaps  IV. Divides the array, sorts the halves and merges them
\optsshort{A-II, B-I, C-III, D-IV}{A-II, B-IV, C-I, D-III}{A-II, B-I, C-IV, D-III}{A-I, B-II, C-IV, D-III}
\end{question}

\begin{question}{M}
Which sorting algorithm makes the minimum number of swaps in the worst case?
\optsshort{Selection sort}{Insertion sort}{Quick sort}{Bubble sort}
\end{question}

\begin{question}{M}
Given below are two statements.\\
Statement I: Merge sort needs O(n) extra space for arrays.\\
Statement II: Merge sort is well suited to sorting linked lists.
\optslong{Both Statement I and Statement II are true}{Both Statement I and Statement II are false}{Statement I is true but Statement II is false}{Statement I is false but Statement II is true}
\end{question}

\begin{question}{M}
Insertion sort is efficient when the input is
\optsshort{random}{nearly sorted}{very large}{reverse sorted}
\end{question}

\begin{question}{M}
The worst-case recursion depth of quick sort on n elements is
\optsshort{O(n)}{O(1)}{O(n log n)}{O(log n)}
\end{question}

\begin{question}{M}
Given below are two statements: one is labelled as Assertion A and the other is labelled as Reason R.\\
Assertion A: Randomized quick sort takes O(n log n) time in the worst case on every input.\\
Reason R: A random pivot makes highly unbalanced partitions unlikely regardless of the input order.
\optslong{Both A and R are true and R is the correct explanation of A}{Both A and R are true but R is NOT the correct explanation of A}{A is true but R is false}{A is false but R is true}
\end{question}

\begin{question}{M}
External sorting is used when
\optslong{the data is already sorted}{the data is very small}{the data has only integers}{the data does not fit in main memory}
\end{question}

\begin{question}{M}
The number of passes of a 2-way external merge sort on 64 initial runs is
\optsshort{8}{6}{7}{5}
\end{question}

\begin{question}{M}
Given below are two statements.\\
Statement I: Heap sort has O(n log n) worst-case time.\\
Statement II: Quick sort is usually faster than heap sort in practice despite its O(n\textsuperscript{2}) worst case.
\optslong{Both Statement I and Statement II are true}{Both Statement I and Statement II are false}{Statement I is true but Statement II is false}{Statement I is false but Statement II is true}
\end{question}

\begin{question}{M}
In the best case, bubble sort with an early-exit flag (stop if no swap in a pass) takes
\optsshort{O(n)}{O(1)}{O(n log n)}{O(n\textsuperscript{2})}
\end{question}

\begin{question}{M}
Linear search on an unsorted array of n elements takes, on average for a successful search,
\optsshort{n\textsuperscript{2} comparisons}{log2 n comparisons}{(n + 1)/2 comparisons}{1 comparison}
\end{question}

\begin{question}{L}
Which sorting algorithm is used by the Unix sort command and many libraries for stable sorting of objects (for example Java's object sort and Python's sort)?
\optslong{Counting sort}{Heap sort}{A merge-sort variant (Timsort)}{Selection sort}
\end{question}

\begin{question}{M}
Given below are two statements.\\
Statement I: Counting sort is a comparison-based algorithm.\\
Statement II: Radix sort can sort in linear time because it does not compare keys.
\optslong{Both Statement I and Statement II are true}{Both Statement I and Statement II are false}{Statement I is true but Statement II is false}{Statement I is false but Statement II is true}
\end{question}

\begin{question}{M}
The number of comparisons made by merge sort on 8 elements already in ascending order (top-down merge sort) is
\optsshort{24}{17}{12}{7}
\end{question}

\begin{question}{M}
Given below are two statements: one is labelled as Assertion A and the other is labelled as Reason R.\\
Assertion A: The worst-case time of any comparison sort on n elements is at least log2(n!).\\
Reason R: log2(n!) = $\Theta$(n log n) by Stirling's approximation.
\optslong{Both A and R are true and R is the correct explanation of A}{Both A and R are true but R is NOT the correct explanation of A}{A is true but R is false}{A is false but R is true}
\end{question}

\begin{question}{M}
Which of the following is an in-place sorting algorithm with O(n log n) worst-case time?
\optsshort{Heap sort}{Counting sort}{Quick sort}{Merge sort}
\end{question}

\begin{question}{M}
For which input does insertion sort perform the maximum number of comparisons?
\optsshort{Sorted ascending}{Sorted descending}{All equal}{Random}
\end{question}

\begin{question}{M}
Sentinel linear search improves performance by
\optslong{sorting the array}{eliminating the index bound check in each iteration}{using binary search}{reducing memory}
\end{question}

\begin{question}{M}
Ternary search on a sorted array of n elements makes about how many comparisons in the worst case compared with binary search?
\optslong{Exactly the same}{More (about 2 log3 n versus log2 n)}{n/3}{Fewer}
\end{question}

\begin{answers}
\ans{1}{2}{Standard complexities of comparison sorts.}
\ans{2}{1}{Quick sort and heap sort can change the relative order of equal keys.}
\ans{3}{2}{A decision tree must have at least n! leaves, so its height is at least log2(n!) = $\Theta$(n log n).}
\ans{4}{3}{$\lceil$log2 5!$\rceil$ = $\lceil$log2 120$\rceil$ = 7, and 7 comparisons are achievable.}
\ans{5}{1}{Each partition splits off only one element.}
\ans{6}{2}{The pivot 4 ends at index 3 with 2, 1, 3 before it and 7, 5, 6, 8 after it.}
\ans{7}{3}{The largest element bubbles to the end: 5 is swapped past 1, 4 and 2.}
\ans{8}{1}{Pass 1 swaps 11 with 64; pass 2 swaps 12 with 25.}
\ans{9}{1}{Count pairs out of order: 7 + 3 + 2 + 4 + 3 + 2 + 1 = 22.}
\ans{10}{4}{Selection sort always makes n(n $-$ 1)/2 comparisons, even on sorted input.}
\ans{11}{4}{A stable sort on the last digit keeps 802 before 2 and 45 before 75.}
\ans{12}{2}{It is linear when k = O(n), and it is stable.}
\ans{13}{1}{It is linear when d is constant and b = O(n).}
\ans{14}{3}{Each bucket then holds O(1) elements on average.}
\ans{15}{3}{$\lfloor$log2 1000$\rfloor$ + 1 = 10.}
\ans{16}{1}{Unsorted data must be sorted first or searched linearly.}
\ans{17}{4}{In the worst case it degrades to O(n).}
\ans{18}{4}{Each comparison outputs one element, and the last element needs no comparison.}
\ans{19}{2}{Process elements in pairs: $\lceil$3n/2$\rceil$ $-$ 2 = 148.}
\ans{20}{3}{A tournament needs n $-$ 1 comparisons for the largest and $\lceil$log2 n$\rceil$ $-$ 1 more: 63 + 5 = 68.}
\ans{21}{3}{Randomized quickselect runs in expected linear time.}
\ans{22}{3}{Basic ideas of common sorting algorithms.}
\ans{23}{1}{It makes at most n $-$ 1 swaps.}
\ans{24}{1}{Linked lists can be merged by relinking nodes, and merge sort needs only sequential access.}
\ans{25}{2}{Its running time is O(n + d), where d is the number of inversions.}
\ans{26}{1}{Recursing on the smaller part first limits the stack depth to O(log n).}
\ans{27}{4}{The O(n log n) bound is on the expected time; with very bad luck the worst case is still O(n\textsuperscript{2}).}
\ans{28}{4}{Sorted runs are created in memory and then merged, often with a k-way merge.}
\ans{29}{2}{64 $\rightarrow$ 32 $\rightarrow$ 16 $\rightarrow$ 8 $\rightarrow$ 4 $\rightarrow$ 2 $\rightarrow$ 1: log2 64 = 6 passes.}
\ans{30}{1}{Quick sort has better cache behaviour and smaller constants.}
\ans{31}{1}{On sorted input, one pass with no swaps suffices.}
\ans{32}{3}{Each position is equally likely.}
\ans{33}{3}{Timsort combines merge sort with insertion sort on runs.}
\ans{34}{4}{Counting sort uses key values as indices rather than comparing them.}
\ans{35}{3}{Merges of sizes 1+1 ($\times$4), 2+2 ($\times$2) and 4+4 ($\times$1) need 4 $\times$ 1 + 2 $\times$ 2 + 1 $\times$ 4 = 12 comparisons.}
\ans{36}{2}{Both are true, but A follows from the decision-tree argument; R only gives the asymptotic size of the bound.}
\ans{37}{1}{Merge sort needs O(n) extra space; quick sort is O(n\textsuperscript{2}) in the worst case.}
\ans{38}{2}{Every new element must be compared with all previous elements.}
\ans{39}{2}{The key is placed at the end so the loop is guaranteed to stop.}
\ans{40}{2}{2 log3 n = 1.26 log2 n, so it is slightly worse.}
\end{answers}

\section{Performance Analysis, Asymptotic Notation and Recurrences}

\begin{question}{H}
Match List I with List II.\\
List I (Notation): A. f(n) = O(g(n))  B. f(n) = $\Omega$(g(n))  C. f(n) = $\Theta$(g(n))  D. f(n) = o(g(n))\\
List II (Meaning): I. g is both an upper and a lower bound for f (up to constants)  II. g is an asymptotic upper bound for f  III. f grows strictly slower than g  IV. g is an asymptotic lower bound for f
\optsshort{A-II, B-I, C-IV, D-III}{A-II, B-IV, C-III, D-I}{A-II, B-IV, C-I, D-III}{A-IV, B-II, C-I, D-III}
\end{question}

\begin{question}{H}
Arrange the following functions in increasing order of growth rate.\\
A. n log n  B. 2\textsuperscript{n}  C. $\surd$n  D. n!  E. log n
\optsshort{C, E, A, B, D}{E, A, C, B, D}{E, C, A, D, B}{E, C, A, B, D}
\end{question}

\begin{question}{H}
Which of the following is true?
\optsshort{n\textsuperscript{2} = O(n log n)}{2\textsuperscript{n+1} = O(2\textsuperscript{n})}{2\textsuperscript{2n} = O(2\textsuperscript{n})}{n! = O(2\textsuperscript{n})}
\end{question}

\begin{question}{H}
log(n!) is
\optsshort{$\Theta$(n)}{$\Theta$(n log n)}{$\Theta$(log n)}{$\Theta$(n\textsuperscript{2})}
\end{question}

\begin{question}{H}
Using the master theorem, the solution of T(n) = 2T(n/2) + n is
\optsshort{$\Theta$(n\textsuperscript{2})}{$\Theta$(log n)}{$\Theta$(n)}{$\Theta$(n log n)}
\end{question}

\begin{question}{H}
Using the master theorem, the solution of T(n) = 4T(n/2) + n is
\optsshort{$\Theta$(n\textsuperscript{2})}{$\Theta$(n)}{$\Theta$(n\textsuperscript{2} log n)}{$\Theta$(n log n)}
\end{question}

\begin{question}{H}
Using the master theorem, the solution of T(n) = 4T(n/2) + n\textsuperscript{2} is
\optsshort{$\Theta$(n log n)}{$\Theta$(n\textsuperscript{2})}{$\Theta$(n\textsuperscript{2} log n)}{$\Theta$(n\textsuperscript{3})}
\end{question}

\begin{question}{H}
Using the master theorem, the solution of T(n) = 4T(n/2) + n\textsuperscript{3} is
\optsshort{$\Theta$(n\textsuperscript{2} log n)}{$\Theta$(n\textsuperscript{3} log n)}{$\Theta$(n\textsuperscript{3})}{$\Theta$(n\textsuperscript{2})}
\end{question}

\begin{question}{H}
The recurrence of Strassen's matrix multiplication, T(n) = 7T(n/2) + $\Theta$(n\textsuperscript{2}), has the solution
\optsshort{$\Theta$(n\textsuperscript{3})}{$\Theta$(n\textsuperscript{2} log n)}{$\Theta$(n\textsuperscript{2})}{$\Theta$(n\textsuperscript{log2 7}) $\approx$ $\Theta$(n\textsuperscript{2}.81)}
\end{question}

\begin{question}{H}
The solution of T(n) = T(n/2) + 1 (binary search) is
\optsshort{$\Theta$(1)}{$\Theta$(n log n)}{$\Theta$(log n)}{$\Theta$(n)}
\end{question}

\begin{question}{H}
The solution of T(n) = T(n $-$ 1) + n, with T(1) = 1, is
\optsshort{$\Theta$(2\textsuperscript{n})}{$\Theta$(n log n)}{$\Theta$(n)}{$\Theta$(n\textsuperscript{2})}
\end{question}

\begin{question}{H}
The solution of T(n) = 2T(n $-$ 1) + 1, with T(1) = 1, is
\optsshort{$\Theta$(n)}{$\Theta$(n!)}{$\Theta$(2\textsuperscript{n})}{$\Theta$(n\textsuperscript{2})}
\end{question}

\begin{question}{H}
The solution of T(n) = 2T($\surd$n) + log n is
\optsshort{$\Theta$(log n)}{$\Theta$(log n log log n)}{$\Theta$($\surd$n)}{$\Theta$(n)}
\end{question}

\begin{question}{M}
The solution of T(n) = T($\surd$n) + 1 is
\optsshort{$\Theta$(log n)}{$\Theta$(log log n)}{$\Theta$(1)}{$\Theta$($\surd$n)}
\end{question}

\begin{question}{M}
The solution of T(n) = 3T(n/4) + n log n is
\optsshort{$\Theta$(n\textsuperscript{log4 3})}{$\Theta$(n)}{$\Theta$(n log\textsuperscript{2} n)}{$\Theta$(n log n)}
\end{question}

\begin{question}{M}
The solution of T(n) = T(n/3) + T(2n/3) + n is
\optsshort{$\Theta$(n log n)}{$\Theta$(n)}{$\Theta$(n\textsuperscript{2})}{$\Theta$(n\textsuperscript{1}.5)}
\end{question}

\begin{question}{M}
The solution of T(n) = 2T(n/2) + n log n is
\optsshort{$\Theta$(n\textsuperscript{2})}{$\Theta$(n log n)}{$\Theta$(n log\textsuperscript{2} n)}{$\Theta$(n log log n)}
\end{question}

\begin{question}{H}
What is the time complexity of the following code?\\
for (i = 1; i < n; i = i * 2) count++;
\optsshort{$\Theta$(n)}{$\Theta$($\surd$n)}{$\Theta$(n log n)}{$\Theta$(log n)}
\end{question}

\begin{question}{H}
What is the time complexity of the following code?\\
for (i = 1; i <= n; i++) for (j = 1; j <= n; j = j + i) count++;
\optsshort{$\Theta$(n)}{$\Theta$(n$\surd$n)}{$\Theta$(n\textsuperscript{2})}{$\Theta$(n log n)}
\end{question}

\begin{question}{H}
What is the time complexity of the following code?\\
for (i = n; i > 0; i = i / 2) for (j = 0; j < i; j++) count++;
\optsshort{$\Theta$(log n)}{$\Theta$(n log n)}{$\Theta$(n\textsuperscript{2})}{$\Theta$(n)}
\end{question}

\begin{question}{H}
What is the time complexity of the following code?\\
for (i = 0; i < n; i++) for (j = 0; j < i; j++) count++;
\optsshort{$\Theta$(n\textsuperscript{3})}{$\Theta$(n\textsuperscript{2})}{$\Theta$(n)}{$\Theta$(n log n)}
\end{question}

\begin{question}{M}
What is the time complexity of the following code?\\
for (i = 1; i * i <= n; i++) count++;
\optsshort{$\Theta$($\surd$n)}{$\Theta$(n\textsuperscript{2})}{$\Theta$(n)}{$\Theta$(log n)}
\end{question}

\begin{question}{M}
The naive recursive Fibonacci function fib(n) = fib(n $-$ 1) + fib(n $-$ 2) runs in
\optslong{$\Theta$(n\textsuperscript{2})}{exponential time $\Theta$($\varphi$\textsuperscript{n}), where $\varphi$ $\approx$ 1.618}{$\Theta$(n)}{$\Theta$(log n)}
\end{question}

\begin{question}{M}
Given below are two statements.\\
Statement I: The space complexity of naive recursive Fibonacci is O(n).\\
Statement II: The time complexity of the iterative (or memoized) Fibonacci is O(n).
\optslong{Both Statement I and Statement II are true}{Both Statement I and Statement II are false}{Statement I is true but Statement II is false}{Statement I is false but Statement II is true}
\end{question}

\begin{question}{H}
Given below are two statements: one is labelled as Assertion A and the other is labelled as Reason R.\\
Assertion A: Appending to a dynamic array that doubles its capacity when full takes O(1) amortized time.\\
Reason R: Every individual append takes O(1) time in the worst case.
\optslong{Both A and R are true and R is the correct explanation of A}{Both A and R are true but R is NOT the correct explanation of A}{A is true but R is false}{A is false but R is true}
\end{question}

\begin{question}{M}
Match List I with List II.\\
List I (Amortized analysis method): A. Aggregate method  B. Accounting method  C. Potential method\\
List II (Idea): I. Assigns a potential function to the data structure state  II. Computes the total cost of n operations and divides by n  III. Charges some operations more than their cost and stores credit for later operations
\optsshort{A-II, B-III, C-I}{A-III, B-II, C-I}{A-I, B-III, C-II}{A-II, B-I, C-III}
\end{question}

\begin{question}{M}
Incrementing a binary counter n times, starting from 0, takes total time
\optsshort{$\Theta$(n\textsuperscript{2})}{$\Theta$(n)}{$\Theta$(log n)}{$\Theta$(n log n)}
\end{question}

\begin{question}{M}
Given below are two statements.\\
Statement I: If f(n) = O(g(n)), then g(n) = O(f(n)).\\
Statement II: If f(n) = O(g(n)) and g(n) = O(h(n)), then f(n) = O(h(n)).
\optslong{Both Statement I and Statement II are true}{Both Statement I and Statement II are false}{Statement I is true but Statement II is false}{Statement I is false but Statement II is true}
\end{question}

\begin{question}{M}
The value of n\textsuperscript{1/log2 n} is
\optsshort{log n}{1}{n}{2}
\end{question}

\begin{question}{M}
Which of the following is NOT O(n\textsuperscript{2})?
\optsshort{n log n}{n}{100n\textsuperscript{2} + 5}{n\textsuperscript{2}.5}
\end{question}

\begin{question}{M}
The master theorem cannot be applied directly to
\optsshort{T(n) = 2\textsuperscript{n} T(n/2) + n}{T(n) = 9T(n/3) + n}{T(n) = 2T(n/2) + n}{T(n) = T(2n/3) + 1}
\end{question}

\begin{question}{M}
Given below are two statements: one is labelled as Assertion A and the other is labelled as Reason R.\\
Assertion A: An algorithm that is O(n\textsuperscript{2}) can never be faster in practice than an O(n log n) algorithm.\\
Reason R: Asymptotic notation ignores constant factors and lower-order terms.
\optslong{Both A and R are true and R is the correct explanation of A}{Both A and R are true but R is NOT the correct explanation of A}{A is true but R is false}{A is false but R is true}
\end{question}

\begin{question}{M}
The time complexity of the following recursive function is\\
void f(int n) \{ if (n <= 1) return; f(n/2); f(n/2); \}
\optsshort{$\Theta$(n\textsuperscript{2})}{$\Theta$(n log n)}{$\Theta$(log n)}{$\Theta$(n)}
\end{question}

\begin{question}{M}
The best-case time complexity of an algorithm is
\optslong{always O(1)}{an upper bound on its running time for every input}{its running time on the input of size n that takes the least time}{the same as its average case}
\end{question}

\begin{question}{M}
Space complexity measures
\optslong{disk usage only}{the number of comparisons}{the total memory an algorithm needs as a function of input size}{the number of instructions}
\end{question}

\begin{question}{M}
Which of the following methods are used to solve recurrences?\\
A. Substitution method  B. Recursion-tree method  C. Master method  D. Greedy method\\
Choose the correct answer from the options given below:
\optsshort{B, C and D only}{A and D only}{A, B, C and D}{A, B and C only}
\end{question}

\begin{question}{M}
The solution of T(n) = 8T(n/2) + n\textsuperscript{2} is
\optsshort{$\Theta$(n\textsuperscript{3})}{$\Theta$(n\textsuperscript{2} log n)}{$\Theta$(n\textsuperscript{2})}{$\Theta$(n\textsuperscript{2}.81)}
\end{question}

\begin{question}{M}
The solution of T(n) = 3T(n/2) + n (Karatsuba multiplication) is
\optsshort{$\Theta$(n log n)}{$\Theta$(n\textsuperscript{log2 3}) $\approx$ $\Theta$(n\textsuperscript{1}.585)}{$\Theta$(n\textsuperscript{1}.5)}{$\Theta$(n\textsuperscript{2})}
\end{question}

\begin{question}{M}
Given below are two statements.\\
Statement I: n log n = O(n\textsuperscript{1}.001).\\
Statement II: (log n)\textsuperscript{100} = O(n\textsuperscript{0}.01).
\optslong{Both Statement I and Statement II are true}{Both Statement I and Statement II are false}{Statement I is true but Statement II is false}{Statement I is false but Statement II is true}
\end{question}

\begin{question}{L}
The term ``asymptotic'' in complexity analysis refers to behaviour as
\optsshort{n grows without bound}{n approaches 0}{n equals 1}{n is negative}
\end{question}

\begin{question}{M}
What is the time complexity of the following code?\\
int i = n; while (i > 1) \{ i = i / 3; \}
\optsshort{$\Theta$($\surd$n)}{$\Theta$(n)}{$\Theta$(log n)}{$\Theta$(n/3)}
\end{question}

\begin{answers}
\ans{1}{3}{Big-O, Big-Omega, Theta and little-o.}
\ans{2}{4}{log n < $\surd$n < n log n < 2\textsuperscript{n} < n!.}
\ans{3}{2}{2\textsuperscript{n+1} = 2 $\times$ 2\textsuperscript{n}, but 2\textsuperscript{2n} = 4\textsuperscript{n} grows faster than 2\textsuperscript{n}.}
\ans{4}{2}{By Stirling's approximation, log(n!) $\approx$ n log n $-$ n log e.}
\ans{5}{4}{n\textsuperscript{log2 2} = n matches f(n), so case 2 gives $\Theta$(n log n).}
\ans{6}{1}{n\textsuperscript{log2 4} = n\textsuperscript{2} dominates f(n) = n (case 1).}
\ans{7}{3}{f(n) = n\textsuperscript{2} equals n\textsuperscript{log2 4}, so case 2 adds a log factor.}
\ans{8}{3}{f(n) = n\textsuperscript{3} dominates n\textsuperscript{2}, and the regularity condition holds (case 3).}
\ans{9}{4}{n\textsuperscript{log2 7} dominates n\textsuperscript{2} (case 1).}
\ans{10}{3}{Case 2 of the master theorem with a = 1, b = 2.}
\ans{11}{4}{T(n) = 1 + 2 + ... + n = n(n + 1)/2.}
\ans{12}{3}{This is the Tower of Hanoi recurrence: T(n) = 2\textsuperscript{n} $-$ 1.}
\ans{13}{2}{Substituting n = 2\textsuperscript{m} gives S(m) = 2S(m/2) + m, so S(m) = $\Theta$(m log m).}
\ans{14}{2}{With n = 2\textsuperscript{m}: S(m) = S(m/2) + 1 = $\Theta$(log m).}
\ans{15}{4}{n\textsuperscript{log4 3} $\approx$ n\textsuperscript{0}.79 is dominated by n log n (case 3).}
\ans{16}{1}{A recursion tree has about log\_(3/2) n levels, each doing O(n) work.}
\ans{17}{3}{This falls between the master theorem cases; the extended case 2 gives $\Theta$(n log\textsuperscript{2} n).}
\ans{18}{4}{i doubles each iteration, so the loop runs about log2 n times.}
\ans{19}{4}{The inner loop runs n/i times: n(1 + 1/2 + ... + 1/n) = $\Theta$(n log n).}
\ans{20}{4}{n + n/2 + n/4 + ... $\leq$ 2n.}
\ans{21}{2}{The total is 0 + 1 + ... + (n $-$ 1) = n(n $-$ 1)/2.}
\ans{22}{1}{The loop stops when i exceeds $\surd$n.}
\ans{23}{2}{Its call tree grows like the Fibonacci numbers themselves.}
\ans{24}{1}{The recursion depth is n even though the number of calls is exponential.}
\ans{25}{3}{An append that triggers resizing copies all elements and takes O(n); only the amortized cost is O(1).}
\ans{26}{1}{The three methods of amortized analysis.}
\ans{27}{2}{Bit i flips n/2\textsuperscript{i} times, so the total is less than 2n.}
\ans{28}{4}{Big-O is transitive but not symmetric (n = O(n\textsuperscript{2}), but n\textsuperscript{2} $\neq$ O(n)).}
\ans{29}{4}{n\textsuperscript{1/log2 n} = 2\textsuperscript{log2 n / log2 n} = 2.}
\ans{30}{4}{n\textsuperscript{2}.5 grows faster than n\textsuperscript{2}.}
\ans{31}{1}{The number of subproblems a must be a constant.}
\ans{32}{4}{For small n or large constant factors an O(n\textsuperscript{2}) algorithm can be faster; R is exactly why A is false.}
\ans{33}{4}{T(n) = 2T(n/2) + 1 = $\Theta$(n).}
\ans{34}{3}{Worst case gives an upper bound over all inputs of size n.}
\ans{35}{3}{Auxiliary space excludes the input itself.}
\ans{36}{4}{The greedy method is an algorithm design technique.}
\ans{37}{1}{n\textsuperscript{log2 8} = n\textsuperscript{3} dominates n\textsuperscript{2} (case 1). This is the recurrence of standard divide-and-conquer matrix multiplication.}
\ans{38}{2}{n\textsuperscript{log2 3} dominates n (case 1).}
\ans{39}{1}{Any power of log n grows more slowly than any positive power of n.}
\ans{40}{1}{Constants and lower-order terms become irrelevant for large n.}
\ans{41}{3}{i is divided by 3 each time: about log3 n iterations.}
\end{answers}

\section{Design Techniques: Divide and Conquer, Greedy Algorithms and Dynamic Programming}

\begin{question}{H}
Match List I with List II.\\
List I (Problem): A. Merge sort  B. Huffman coding  C. Matrix-chain multiplication  D. N-queens\\
List II (Design technique): I. Dynamic programming  II. Greedy method  III. Divide and conquer  IV. Backtracking
\optsshort{A-III, B-II, C-IV, D-I}{A-III, B-I, C-II, D-IV}{A-II, B-III, C-I, D-IV}{A-III, B-II, C-I, D-IV}
\end{question}

\begin{question}{H}
Strassen's algorithm multiplies two 2 $\times$ 2 block matrices using
\optsshort{8 multiplications}{7 multiplications}{4 multiplications}{6 multiplications}
\end{question}

\begin{question}{M}
Karatsuba's algorithm multiplies two n-digit numbers using
\optslong{three half-size multiplications, giving $\Theta$(n\textsuperscript{1}.585)}{two half-size multiplications}{four half-size multiplications}{n multiplications}
\end{question}

\begin{question}{H}
Given below are two statements.\\
Statement I: Dynamic programming is applicable when a problem has optimal substructure and overlapping subproblems.\\
Statement II: Divide and conquer is most efficient when the subproblems overlap heavily.
\optslong{Both Statement I and Statement II are true}{Both Statement I and Statement II are false}{Statement I is true but Statement II is false}{Statement I is false but Statement II is true}
\end{question}

\begin{question}{H}
The minimum number of scalar multiplications needed to compute the product of matrices with dimensions 10 $\times$ 100, 100 $\times$ 5 and 5 $\times$ 50 is
\optsshort{25,000}{5,000}{75,000}{7,500}
\end{question}

\begin{question}{H}
For the matrix chain with dimensions p = (30, 35, 15, 5, 10, 20, 25), the minimum number of scalar multiplications is
\optsshort{18,000}{15,125}{15,000}{11,875}
\end{question}

\begin{question}{M}
For three matrices of dimensions 10 $\times$ 20, 20 $\times$ 30 and 30 $\times$ 40, the minimum number of scalar multiplications is
\optsshort{30,000}{24,000}{18,000}{32,000}
\end{question}

\begin{question}{M}
The number of ways to fully parenthesize a product of 4 matrices is
\optsshort{4}{6}{5}{3}
\end{question}

\begin{question}{H}
The length of the longest common subsequence of ``ABCBDAB'' and ``BDCABA'' is
\optsshort{6}{5}{4}{3}
\end{question}

\begin{question}{M}
The length of the longest common subsequence of ``AGGTAB'' and ``GXTXAYB'' is
\optsshort{6}{3}{4}{5}
\end{question}

\begin{question}{H}
The time complexity of the dynamic-programming algorithm for the LCS of strings of lengths m and n is
\optsshort{O(m + n)}{O(mn)}{O(m log n)}{O(2\textsuperscript{m+n})}
\end{question}

\begin{question}{H}
The edit (Levenshtein) distance between ``kitten'' and ``sitting'' is
\optsshort{2}{4}{5}{3}
\end{question}

\begin{question}{H}
For the 0/1 knapsack with capacity 50 and items (weight, value) = (10, 60), (20, 100), (30, 120), the maximum value is
\optsshort{220}{240}{160}{180}
\end{question}

\begin{question}{H}
For the fractional knapsack with capacity 50 and items (weight, value) = (10, 60), (20, 100), (30, 120), the maximum value is
\optsshort{220}{280}{230}{240}
\end{question}

\begin{question}{M}
For the 0/1 knapsack with capacity 7 and items (weight, value) = (1, 1), (3, 4), (4, 5), (5, 7), the maximum value is
\optsshort{10}{11}{9}{8}
\end{question}

\begin{question}{H}
Given below are two statements: one is labelled as Assertion A and the other is labelled as Reason R.\\
Assertion A: The greedy method of choosing items by value-to-weight ratio is optimal for the fractional knapsack problem but not for the 0/1 knapsack problem.\\
Reason R: In the 0/1 knapsack, an item cannot be split, so filling the remaining capacity exactly may require a different selection.
\optslong{Both A and R are true and R is the correct explanation of A}{Both A and R are true but R is NOT the correct explanation of A}{A is true but R is false}{A is false but R is true}
\end{question}

\begin{question}{H}
The dynamic-programming algorithm for 0/1 knapsack with n items and capacity W runs in
\optslong{O(n log n)}{O(n + W)}{O(nW), which is pseudo-polynomial}{O(2\textsuperscript{n}) always}
\end{question}

\begin{question}{H}
Huffman coding for characters with frequencies a: 45, b: 13, c: 12, d: 16, e: 9, f: 5 (per 100 characters) gives an average code length of
\optsshort{2.5 bits}{2.24 bits}{3 bits}{2.0 bits}
\end{question}

\begin{question}{M}
Huffman coding for symbols with probabilities 0.4, 0.2, 0.2, 0.1, 0.1 gives an average code length of
\optsshort{2.4 bits}{2.2 bits}{2.0 bits}{3.0 bits}
\end{question}

\begin{question}{H}
In the activity-selection problem with activities (start, finish): (1, 4), (3, 5), (0, 6), (5, 7), (3, 9), (5, 9), (6, 10), (8, 11), (8, 12), (2, 14), (12, 16), the maximum number of compatible activities is
\optsshort{3}{5}{6}{4}
\end{question}

\begin{question}{H}
Jobs J1 to J5 have (deadline, profit): (2, 100), (1, 19), (2, 27), (1, 25), (3, 15), and each takes one unit of time. The maximum total profit using greedy job sequencing is
\optsshort{140}{152}{142}{167}
\end{question}

\begin{question}{M}
Using the optimal merge pattern, sorted files of sizes 20, 30, 10, 5 and 30 records are merged two at a time. The minimum total number of record moves is
\optsshort{95}{165}{240}{205}
\end{question}

\begin{question}{H}
With coin denominations \{1, 3, 4\} and amount 6, which statement is correct?
\optslong{Greedy gives 2 coins, which is optimal}{Greedy gives 3 coins (4 + 1 + 1), but the optimum is 2 coins (3 + 3)}{Both give 3 coins}{No solution exists}
\end{question}

\begin{question}{M}
The number of ways to make change for 5 using coins \{1, 2, 5\} (unlimited supply, order ignored) is
\optsshort{6}{5}{4}{3}
\end{question}

\begin{question}{H}
Match List I with List II.\\
List I (Algorithm): A. Floyd-Warshall  B. Bellman-Ford  C. Kruskal  D. Dijkstra\\
List II (Technique): I. Greedy, single-source shortest paths with non-negative weights  II. Dynamic programming, single-source shortest paths with negative weights allowed  III. Dynamic programming, all-pairs shortest paths  IV. Greedy, minimum spanning tree
\optsshort{A-III, B-IV, C-II, D-I}{A-II, B-III, C-IV, D-I}{A-III, B-II, C-IV, D-I}{A-III, B-II, C-I, D-IV}
\end{question}

\begin{question}{M}
The dynamic-programming solution of the travelling salesman problem (Held-Karp) runs in
\optsshort{O(n\textsuperscript{3})}{O(n!)}{O(n\textsuperscript{2} 2\textsuperscript{n})}{O(2\textsuperscript{n})}
\end{question}

\begin{question}{M}
The dynamic-programming algorithm for constructing an optimal binary search tree on n keys runs in
\optslong{O(n\textsuperscript{3}), or O(n\textsuperscript{2}) with Knuth's optimization}{O(2\textsuperscript{n})}{O(n\textsuperscript{2})}{O(n log n)}
\end{question}

\begin{question}{M}
The longest increasing subsequence of a sequence of n numbers can be found in
\optsshort{O(n\textsuperscript{1}.5)}{O(n log n)}{O(n\textsuperscript{3})}{O(2\textsuperscript{n})}
\end{question}

\begin{question}{M}
Given below are two statements.\\
Statement I: Memoization is a top-down approach that stores the results of recursive calls.\\
Statement II: Tabulation is a bottom-up approach that fills a table iteratively.
\optslong{Both Statement I and Statement II are true}{Both Statement I and Statement II are false}{Statement I is true but Statement II is false}{Statement I is false but Statement II is true}
\end{question}

\begin{question}{M}
The principle of optimality states that
\optslong{an optimal solution contains optimal solutions to its subproblems}{a greedy choice is always optimal}{the optimal solution is unique}{every problem can be solved by divide and conquer}
\end{question}

\begin{question}{M}
Given below are two statements: one is labelled as Assertion A and the other is labelled as Reason R.\\
Assertion A: A greedy algorithm always produces an optimal solution.\\
Reason R: A greedy algorithm makes the locally best choice at each step.
\optslong{Both A and R are true and R is the correct explanation of A}{Both A and R are true but R is NOT the correct explanation of A}{A is true but R is false}{A is false but R is true}
\end{question}

\begin{question}{M}
The maximum subarray problem can be solved in linear time by
\optsshort{Huffman coding}{binary search}{Kadane's algorithm}{merge sort}
\end{question}

\begin{question}{M}
The closest pair of points among n points in the plane can be found by divide and conquer in
\optsshort{O(n\textsuperscript{3})}{O(n log n)}{O(n)}{O(n\textsuperscript{2})}
\end{question}

\begin{question}{M}
Given below are two statements.\\
Statement I: Dijkstra's algorithm uses dynamic programming with backtracking.\\
Statement II: Prim's and Kruskal's algorithms are both greedy.
\optslong{Both Statement I and Statement II are true}{Both Statement I and Statement II are false}{Statement I is true but Statement II is false}{Statement I is false but Statement II is true}
\end{question}

\begin{question}{M}
The number of subproblems solved by the dynamic-programming algorithm for matrix-chain multiplication of n matrices is
\optsshort{O(2\textsuperscript{n})}{O(n)}{O(n\textsuperscript{2})}{O(n\textsuperscript{3})}
\end{question}

\begin{question}{M}
The edit distance between ``sunday'' and ``saturday'' is
\optsshort{4}{5}{3}{2}
\end{question}

\begin{question}{M}
Given below are two statements: one is labelled as Assertion A and the other is labelled as Reason R.\\
Assertion A: Huffman coding produces a prefix code.\\
Reason R: Huffman's algorithm repeatedly merges the two least frequent nodes using a priority queue.
\optslong{Both A and R are true and R is the correct explanation of A}{Both A and R are true but R is NOT the correct explanation of A}{A is true but R is false}{A is false but R is true}
\end{question}

\begin{question}{L}
Binary search is an example of which design technique?
\optslong{Backtracking}{Divide and conquer (decrease and conquer)}{Greedy}{Dynamic programming}
\end{question}

\begin{answers}
\ans{1}{4}{Standard classification of algorithms by design technique.}
\ans{2}{2}{It trades one multiplication for extra additions, giving $\Theta$(n\textsuperscript{2}.81).}
\ans{3}{1}{T(n) = 3T(n/2) + O(n).}
\ans{4}{3}{Divide and conquer suits independent subproblems; overlapping subproblems are what DP exploits.}
\ans{5}{4}{(A1 A2) A3 costs 10$\cdot$100$\cdot$5 + 10$\cdot$5$\cdot$50 = 5000 + 2500 = 7500; A1 (A2 A3) costs 75,000.}
\ans{6}{2}{This is the CLRS example; the optimal parenthesization is ((A1 (A2 A3)) ((A4 A5) A6)).}
\ans{7}{3}{(A1 A2) A3: 6000 + 12000 = 18,000; A1 (A2 A3): 24000 + 8000 = 32,000.}
\ans{8}{3}{It is the Catalan number C(n $-$ 1) = C(3) = 5.}
\ans{9}{3}{One LCS is BCBA.}
\ans{10}{3}{GTAB is a common subsequence of length 4.}
\ans{11}{2}{It fills an (m + 1) $\times$ (n + 1) table.}
\ans{12}{4}{Substitute k $\rightarrow$ s, substitute e $\rightarrow$ i, and insert g.}
\ans{13}{1}{Take the items of weight 20 and 30; 240 is the fractional-knapsack optimum.}
\ans{14}{4}{Take items 1 and 2 fully (value 160) and 20/30 of item 3 (value 80).}
\ans{15}{3}{Items of weight 3 and 4 give value 9.}
\ans{16}{1}{0/1 knapsack needs dynamic programming (O(nW), pseudo-polynomial).}
\ans{17}{3}{It is exponential in the number of bits of W.}
\ans{18}{2}{Code lengths are a: 1, b, c, d: 3, e, f: 4; total 224 bits per 100 characters.}
\ans{19}{2}{One optimal set of lengths is 2, 2, 2, 3, 3: 0.8 + 0.4 + 0.4 + 0.3 + 0.3 = 2.2.}
\ans{20}{4}{Greedy by earliest finish: (1, 4), (5, 7), (8, 11), (12, 16).}
\ans{21}{3}{Schedule J3 in slot 1, J1 in slot 2 and J5 in slot 3: 27 + 100 + 15 = 142.}
\ans{22}{4}{Merge 5 + 10 = 15, 15 + 20 = 35, 30 + 30 = 60, 35 + 60 = 95: 15 + 35 + 60 + 95 = 205.}
\ans{23}{2}{Greedy coin change is optimal only for special (canonical) coin systems.}
\ans{24}{3}{5; 2 + 2 + 1; 2 + 1 + 1 + 1; 1 + 1 + 1 + 1 + 1.}
\ans{25}{3}{Techniques behind common graph algorithms.}
\ans{26}{3}{It is still exponential but much faster than trying all n! tours.}
\ans{27}{1}{It fills a table of costs for every range of keys.}
\ans{28}{2}{The simple DP takes O(n\textsuperscript{2}); patience sorting with binary search takes O(n log n).}
\ans{29}{1}{Both are ways of implementing dynamic programming.}
\ans{30}{1}{It is the basis of dynamic programming (Bellman).}
\ans{31}{4}{Greedy algorithms are optimal only for problems with the greedy-choice property and optimal substructure.}
\ans{32}{3}{The divide-and-conquer solution takes O(n log n).}
\ans{33}{2}{The brute-force method takes O(n\textsuperscript{2}).}
\ans{34}{4}{Dijkstra's algorithm is a greedy algorithm.}
\ans{35}{3}{One subproblem per pair (i, j) with i $\leq$ j, each solved in O(n) time, giving O(n\textsuperscript{3}) total.}
\ans{36}{3}{Insert a, insert t, substitute n $\rightarrow$ r.}
\ans{37}{2}{Both are true, but the prefix property comes from characters being stored only at the leaves of the tree.}
\ans{38}{2}{It reduces the problem to one subproblem of half the size.}
\end{answers}

\section{Backtracking, Branch and Bound, and Lower Bound Theory}

\begin{question}{H}
The number of distinct solutions of the 4-queens problem is
\optsshort{2}{4}{1}{8}
\end{question}

\begin{question}{H}
The number of distinct solutions of the 8-queens problem is
\optsshort{12}{92}{64}{96}
\end{question}

\begin{question}{H}
For the sum-of-subsets problem with w = \{5, 10, 12, 13, 15, 18\} and target M = 30, the number of solutions is
\optsshort{4}{5}{2}{3}
\end{question}

\begin{question}{H}
Given below are two statements.\\
Statement I: Backtracking explores the state-space tree in depth-first order and abandons a partial solution as soon as it cannot lead to a valid solution.\\
Statement II: Branch and bound typically explores the state-space tree using breadth-first or least-cost search and uses bounds to prune nodes.
\optslong{Both Statement I and Statement II are true}{Both Statement I and Statement II are false}{Statement I is true but Statement II is false}{Statement I is false but Statement II is true}
\end{question}

\begin{question}{H}
Match List I with List II.\\
List I (Branch-and-bound strategy): A. FIFO branch and bound  B. LIFO branch and bound  C. Least-cost (LC) branch and bound\\
List II (Order in which live nodes are expanded): I. The node with the best estimated cost is expanded next  II. Live nodes are kept in a queue  III. Live nodes are kept in a stack
\optsshort{A-II, B-I, C-III}{A-III, B-II, C-I}{A-I, B-III, C-II}{A-II, B-III, C-I}
\end{question}

\begin{question}{M}
In backtracking, a function that decides whether a node of the state-space tree can possibly lead to a solution is called a
\optsshort{objective function}{cost function}{bounding function}{hash function}
\end{question}

\begin{question}{M}
Which of the following problems is commonly solved by backtracking?\\
A. Graph m-colouring  B. Hamiltonian cycle  C. Sum of subsets  D. Minimum spanning tree\\
Choose the correct answer from the options given below:
\optsshort{A and D only}{A, B, C and D}{A, B and C only}{B, C and D only}
\end{question}

\begin{question}{M}
In the m-colouring decision problem solved by backtracking, a vertex is given a colour only if
\optslong{it has the highest degree}{no adjacent vertex already has that colour}{the graph is a tree}{all vertices are uncoloured}
\end{question}

\begin{question}{M}
In the state-space tree of the n-queens problem with one queen per row, the number of leaf nodes before pruning is
\optslong{n}{n\textsuperscript{2}}{n\textsuperscript{n} (or n! if each column is used once)}{2\textsuperscript{n}}
\end{question}

\begin{question}{H}
Given below are two statements: one is labelled as Assertion A and the other is labelled as Reason R.\\
Assertion A: Branch and bound is suited to optimization problems such as the travelling salesman problem.\\
Reason R: Branch and bound compares the bound at each node with the best solution found so far and discards nodes that cannot do better.
\optslong{Both A and R are true and R is the correct explanation of A}{Both A and R are true but R is NOT the correct explanation of A}{A is true but R is false}{A is false but R is true}
\end{question}

\begin{question}{M}
The 15-puzzle is typically solved using
\optslong{Huffman coding}{Kruskal's algorithm}{LC branch and bound with a cost estimate such as the number of misplaced tiles}{merge sort}
\end{question}

\begin{question}{H}
The decision tree for any comparison-based algorithm that sorts n elements must have at least
\optsshort{2\textsuperscript{n} leaves}{n\textsuperscript{2} leaves}{n! leaves}{n leaves}
\end{question}

\begin{question}{H}
The lower bound on the worst-case number of comparisons for searching a sorted array of n elements by comparisons is
\optsshort{$\Omega$(1)}{$\Omega$(n)}{$\Omega$(n log n)}{$\Omega$(log n)}
\end{question}

\begin{question}{H}
Finding the maximum of n elements by comparisons requires at least
\optsshort{log2 n comparisons}{n log n comparisons}{n/2 comparisons}{n $-$ 1 comparisons}
\end{question}

\begin{question}{M}
The minimum height of a comparison tree that sorts 8 elements is
\optsshort{24}{15}{28}{16}
\end{question}

\begin{question}{M}
Merging two sorted lists of n elements each needs, in the worst case, at least
\optsshort{n log n comparisons}{log n comparisons}{n comparisons}{2n $-$ 1 comparisons}
\end{question}

\begin{question}{H}
Given below are two statements.\\
Statement I: The element-uniqueness problem has an $\Omega$(n log n) lower bound in the algebraic decision-tree model.\\
Statement II: Computing the convex hull of n points has an $\Omega$(n log n) lower bound, shown by reduction from sorting.
\optslong{Both Statement I and Statement II are true}{Both Statement I and Statement II are false}{Statement I is true but Statement II is false}{Statement I is false but Statement II is true}
\end{question}

\begin{question}{M}
Proving a lower bound for problem B by reduction means showing that
\optslong{B is in P}{B can be solved quickly}{B has no solution}{a problem A with a known lower bound can be transformed into B efficiently, so B must be at least as hard}
\end{question}

\begin{question}{M}
In an adversary argument for a lower bound, the adversary
\optslong{always gives random answers}{solves the problem}{answers the algorithm's queries so as to force it to do as much work as possible, while staying consistent}{chooses the input in advance}
\end{question}

\begin{question}{M}
Given below are two statements: one is labelled as Assertion A and the other is labelled as Reason R.\\
Assertion A: Merge sort is asymptotically optimal among comparison sorts.\\
Reason R: Merge sort uses O(n) extra memory.
\optslong{Both A and R are true and R is the correct explanation of A}{Both A and R are true but R is NOT the correct explanation of A}{A is true but R is false}{A is false but R is true}
\end{question}

\begin{question}{M}
Counting sort beats the $\Omega$(n log n) lower bound because
\optslong{it does not compare keys, using them as array indices instead}{it uses recursion}{its input is always small}{it is a comparison sort}
\end{question}

\begin{question}{M}
Hamiltonian-cycle backtracking adds a vertex to the path only if
\optslong{it is adjacent to the last vertex of the path and not yet on the path}{it has the smallest label}{it has been visited before}{it is a leaf}
\end{question}

\begin{question}{M}
In the 0/1 knapsack branch-and-bound method, the upper bound at a node is usually computed by
\optslong{solving the fractional knapsack on the remaining items}{the value of items already chosen}{counting the remaining items}{dynamic programming on all items}
\end{question}

\begin{question}{M}
Given below are two statements.\\
Statement I: Backtracking always runs in polynomial time.\\
Statement II: Branch and bound always runs in polynomial time.
\optslong{Both Statement I and Statement II are true}{Both Statement I and Statement II are false}{Statement I is true but Statement II is false}{Statement I is false but Statement II is true}
\end{question}

\begin{question}{L}
The term ``state-space tree'' refers to
\optslong{a tree representing all possible partial and complete solutions examined by the algorithm}{the call stack}{a binary search tree}{a heap}
\end{question}

\begin{question}{M}
In a game tree, alpha-beta pruning
\optslong{explores every node}{skips branches that cannot affect the final minimax decision}{only works for single-player games}{changes the minimax value}
\end{question}

\begin{question}{M}
The graph m-colouring problem for a graph with n vertices solved by backtracking has a worst-case time of
\optsshort{O(n\textsuperscript{2})}{O(nm)}{O(m\textsuperscript{n} $\times$ n)}{O(n log m)}
\end{question}

\begin{answers}
\ans{1}{1}{The two solutions are mirror images: (2, 4, 1, 3) and (3, 1, 4, 2).}
\ans{2}{2}{There are 92 solutions, of which 12 are distinct up to rotation and reflection.}
\ans{3}{4}{\{5, 10, 15\}, \{5, 12, 13\} and \{12, 18\}.}
\ans{4}{1}{Both are standard descriptions of the two techniques.}
\ans{5}{4}{LC search uses a priority queue ordered by a cost function.}
\ans{6}{3}{Nodes failing the bounding function are pruned.}
\ans{7}{3}{The minimum spanning tree has efficient greedy algorithms.}
\ans{8}{2}{The algorithm backtracks when no colour is valid for the current vertex.}
\ans{9}{3}{Backtracking prunes most of these.}
\ans{10}{1}{Backtracking, in contrast, is mainly used for decision and enumeration problems.}
\ans{11}{3}{Not all initial arrangements are solvable; a parity test decides solvability.}
\ans{12}{3}{Each permutation of the input must lead to a different leaf.}
\ans{13}{4}{The decision tree must have at least n + 1 leaves.}
\ans{14}{4}{Every element except the maximum must lose at least one comparison.}
\ans{15}{4}{$\lceil$log2 8!$\rceil$ = $\lceil$log2 40320$\rceil$ = $\lceil$15.3$\rceil$ = 16.}
\ans{16}{4}{An adversary argument shows that every adjacent pair in the output must be compared.}
\ans{17}{1}{If convex hull could be found faster, sorting could be done faster by mapping x to (x, x\textsuperscript{2}).}
\ans{18}{4}{If B were faster, A could be solved faster too, contradicting A's lower bound.}
\ans{19}{3}{The adversary builds a worst-case input on the fly.}
\ans{20}{2}{Both are true, but A holds because merge sort's O(n log n) matches the $\Omega$(n log n) lower bound; its memory use is unrelated.}
\ans{21}{1}{The decision-tree lower bound applies only to comparison sorts.}
\ans{22}{1}{At the end the last vertex must also be adjacent to the start.}
\ans{23}{1}{The fractional relaxation gives an optimistic bound.}
\ans{24}{2}{Both can take exponential time in the worst case; pruning helps only in practice.}
\ans{25}{1}{Its nodes are problem states and its edges are choices.}
\ans{26}{2}{With best move ordering it searches about O(b\textsuperscript{d/2}) nodes.}
\ans{27}{3}{There are m\textsuperscript{n} possible colourings, each checked in O(n) or more.}
\end{answers}

\section{Graph Algorithms (representations, BFS, DFS, topological sort, MST, shortest paths, maximum flow)}

\begin{question}{H}
Match List I with List II.\\
List I (Graph algorithm): A. Breadth-first search  B. Dijkstra's algorithm  C. Bellman-Ford algorithm  D. Floyd-Warshall algorithm\\
List II (Time complexity): I. O(V\textsuperscript{3})  II. O(V + E)  III. O(VE)  IV. O((V + E) log V) with a binary heap
\optsshort{A-IV, B-II, C-III, D-I}{A-II, B-IV, C-III, D-I}{A-II, B-IV, C-I, D-III}{A-II, B-III, C-IV, D-I}
\end{question}

\begin{question}{H}
For a graph with V vertices and E edges, the space needed by an adjacency matrix and by adjacency lists is
\optsshort{O(V\textsuperscript{2}) and O(V + E)}{O(E) and O(V)}{O(V + E) and O(V\textsuperscript{2})}{O(V log V) and O(E)}
\end{question}

\begin{question}{H}
The minimum spanning tree of the CLRS example graph with edges a-b 4, a-h 8, b-c 8, b-h 11, c-d 7, c-f 4, c-i 2, d-e 9, d-f 14, e-f 10, f-g 2, g-h 1, g-i 6, h-i 7 has total weight
\optsshort{35}{41}{37}{39}
\end{question}

\begin{question}{M}
A graph has edges A-B 1, A-C 3, B-C 1, B-D 4, C-D 2, C-E 5, D-E 6, D-F 7, E-F 3. The weight of its minimum spanning tree is
\optsshort{13}{11}{12}{14}
\end{question}

\begin{question}{H}
Given below are two statements.\\
Statement I: If all edge weights of a connected graph are distinct, its minimum spanning tree is unique.\\
Statement II: A minimum spanning tree of a graph with V vertices has V edges.
\optslong{Both Statement I and Statement II are true}{Both Statement I and Statement II are false}{Statement I is true but Statement II is false}{Statement I is false but Statement II is true}
\end{question}

\begin{question}{H}
Match List I with List II.\\
List I (MST algorithm): A. Kruskal  B. Prim with a binary heap  C. Prim with a Fibonacci heap\\
List II (Time): I. O(E + V log V)  II. O(E log E)  III. O(E log V)
\optsshort{A-I, B-III, C-II}{A-III, B-II, C-I}{A-II, B-I, C-III}{A-II, B-III, C-I}
\end{question}

\begin{question}{H}
In the CLRS graph with edges s$\rightarrow$t 10, s$\rightarrow$y 5, t$\rightarrow$x 1, t$\rightarrow$y 2, y$\rightarrow$t 3, y$\rightarrow$x 9, y$\rightarrow$z 2, x$\rightarrow$z 4, z$\rightarrow$s 7, z$\rightarrow$x 6, Dijkstra's algorithm from s gives the shortest distances to t, x, y, z as
\optsshort{8, 9, 5, 7}{10, 11, 5, 7}{8, 13, 5, 7}{8, 9, 5, 6}
\end{question}

\begin{question}{M}
In the CLRS graph with edges s$\rightarrow$t 10, s$\rightarrow$y 5, t$\rightarrow$x 1, t$\rightarrow$y 2, y$\rightarrow$t 3, y$\rightarrow$x 9, y$\rightarrow$z 2, x$\rightarrow$z 4, z$\rightarrow$s 7, z$\rightarrow$x 6, the order in which Dijkstra's algorithm from s finalizes the vertices is
\optsshort{s, y, z, t, x}{s, y, t, x, z}{s, t, y, x, z}{s, y, t, z, x}
\end{question}

\begin{question}{H}
Given below are two statements: one is labelled as Assertion A and the other is labelled as Reason R.\\
Assertion A: Dijkstra's algorithm may give wrong results on graphs with negative edge weights.\\
Reason R: Once a vertex is finalized, Dijkstra's algorithm never updates its distance, but a later negative edge could provide a shorter path.
\optslong{Both A and R are true and R is the correct explanation of A}{Both A and R are true but R is NOT the correct explanation of A}{A is true but R is false}{A is false but R is true}
\end{question}

\begin{question}{H}
Given below are two statements.\\
Statement I: The Bellman-Ford algorithm can detect negative-weight cycles reachable from the source.\\
Statement II: Bellman-Ford relaxes every edge V $-$ 1 times.
\optslong{Both Statement I and Statement II are true}{Both Statement I and Statement II are false}{Statement I is true but Statement II is false}{Statement I is false but Statement II is true}
\end{question}

\begin{question}{M}
Running Bellman-Ford from s on the CLRS graph with edges s$\rightarrow$t 6, s$\rightarrow$y 7, t$\rightarrow$x 5, t$\rightarrow$y 8, t$\rightarrow$z $-$4, x$\rightarrow$t $-$2, y$\rightarrow$x $-$3, y$\rightarrow$z 9, z$\rightarrow$s 2, z$\rightarrow$x 7 gives the shortest distance to z as
\optsshort{2}{$-$4}{4}{$-$2}
\end{question}

\begin{question}{H}
BFS on an unweighted graph from a source s finds
\optslong{a minimum spanning tree}{all strongly connected components}{the shortest path (fewest edges) from s to every reachable vertex}{a topological order}
\end{question}

\begin{question}{H}
For the undirected graph with adjacency lists A: B, C; B: A, D, E; C: A, F; D: B; E: B, F; F: C, E, the BFS and DFS orders starting from A (taking neighbours in list order) are
\optslong{BFS: A B C D E F; DFS: A B C F E D}{BFS: A B C D E F; DFS: A B D E F C}{BFS: A C B F E D; DFS: A B D E F C}{BFS: A B D E F C; DFS: A B C D E F}
\end{question}

\begin{question}{H}
Given below are two statements.\\
Statement I: DFS of an undirected graph can produce cross edges.\\
Statement II: A directed graph has a cycle if and only if DFS finds a back edge.
\optslong{Both Statement I and Statement II are true}{Both Statement I and Statement II are false}{Statement I is true but Statement II is false}{Statement I is false but Statement II is true}
\end{question}

\begin{question}{H}
The number of topological orderings of the DAG with edges A$\rightarrow$B, A$\rightarrow$C, B$\rightarrow$D, C$\rightarrow$D, D$\rightarrow$E is
\optsshort{2}{3}{1}{4}
\end{question}

\begin{question}{M}
The number of topological orderings of the DAG with edges 1$\rightarrow$2, 1$\rightarrow$3, 2$\rightarrow$4, 3$\rightarrow$4, 4$\rightarrow$5, 4$\rightarrow$6 is
\optsshort{8}{4}{6}{2}
\end{question}

\begin{question}{H}
A topological sort exists for a directed graph if and only if the graph is
\optsshort{strongly connected}{complete}{acyclic}{bipartite}
\end{question}

\begin{question}{M}
Kosaraju's algorithm finds the strongly connected components of a directed graph using
\optslong{one BFS}{union-find}{two DFS passes, the second on the transpose graph in decreasing order of finishing time}{Dijkstra's algorithm}
\end{question}

\begin{question}{H}
The value of the maximum flow in the CLRS network with edges s$\rightarrow$v1 16, s$\rightarrow$v2 13, v1$\rightarrow$v3 12, v2$\rightarrow$v1 4, v2$\rightarrow$v4 14, v3$\rightarrow$v2 9, v3$\rightarrow$t 20, v4$\rightarrow$v3 7, v4$\rightarrow$t 4 is
\optsshort{26}{19}{23}{29}
\end{question}

\begin{question}{M}
A network has edges S$\rightarrow$A 10, S$\rightarrow$B 5, A$\rightarrow$B 15, A$\rightarrow$T 5, B$\rightarrow$T 10. The maximum flow from S to T is
\optsshort{25}{20}{10}{15}
\end{question}

\begin{question}{H}
Given below are two statements: one is labelled as Assertion A and the other is labelled as Reason R.\\
Assertion A: The value of a maximum flow equals the capacity of a minimum cut.\\
Reason R: When no augmenting path remains in the residual network, the vertices reachable from the source define a cut whose capacity equals the flow.
\optslong{Both A and R are true and R is the correct explanation of A}{Both A and R are true but R is NOT the correct explanation of A}{A is true but R is false}{A is false but R is true}
\end{question}

\begin{question}{M}
The Edmonds-Karp algorithm, which always augments along a shortest path found by BFS, runs in
\optsshort{O(E log V)}{O(VE\textsuperscript{2})}{O(V\textsuperscript{3})}{O(E |f*|)}
\end{question}

\begin{question}{M}
Maximum bipartite matching can be solved by
\optslong{Huffman coding}{reducing it to a maximum-flow problem}{Dijkstra's algorithm}{topological sorting}
\end{question}

\begin{question}{M}
The Floyd-Warshall algorithm on the graph with distance matrix rows [0, 3, $\infty$, 7], [8, 0, 2, $\infty$], [5, $\infty$, 0, 1], [2, $\infty$, $\infty$, 0] gives the shortest distance from vertex 1 to vertex 4 as
\optsshort{6}{7}{8}{5}
\end{question}

\begin{question}{H}
The number of edges in a complete undirected graph with 10 vertices is
\optsshort{45}{20}{90}{100}
\end{question}

\begin{question}{M}
By Cayley's formula, the number of labelled spanning trees of the complete graph K5 is
\optsshort{25}{120}{60}{125}
\end{question}

\begin{question}{H}
Given below are two statements.\\
Statement I: A connected undirected graph has an Euler circuit if and only if every vertex has even degree.\\
Statement II: Determining whether a graph has a Hamiltonian cycle can be done in polynomial time.
\optslong{Both Statement I and Statement II are true}{Both Statement I and Statement II are false}{Statement I is true but Statement II is false}{Statement I is false but Statement II is true}
\end{question}

\begin{question}{M}
An undirected graph is bipartite if and only if it
\optslong{has an even number of vertices}{is complete}{contains no cycle of odd length}{is a tree}
\end{question}

\begin{question}{M}
The sum of the degrees of all vertices of an undirected graph with E edges is
\optsshort{E}{2E}{V + E}{E\textsuperscript{2}}
\end{question}

\begin{question}{M}
An articulation point (cut vertex) of a connected undirected graph is a vertex whose removal
\optsshort{disconnects the graph}{creates a cycle}{makes the graph complete}{has no effect}
\end{question}

\begin{question}{M}
Johnson's algorithm for all-pairs shortest paths in sparse graphs with negative edges
\optslong{reweights edges using Bellman-Ford and then runs Dijkstra from every vertex}{uses only BFS}{cannot handle negative edges}{uses Floyd-Warshall}
\end{question}

\begin{question}{M}
Given below are two statements.\\
Statement I: Prim's algorithm grows a single tree from a starting vertex.\\
Statement II: Kruskal's algorithm may maintain a forest of several trees during its execution.
\optslong{Both Statement I and Statement II are true}{Both Statement I and Statement II are false}{Statement I is true but Statement II is false}{Statement I is false but Statement II is true}
\end{question}

\begin{question}{M}
Warshall's algorithm computes
\optslong{shortest paths}{a minimum spanning tree}{maximum flow}{the transitive closure of a directed graph in O(V\textsuperscript{3})}
\end{question}

\begin{question}{M}
Which data structure is used by Dijkstra's algorithm to select the next vertex efficiently?
\optsshort{Hash table}{Min-priority queue}{Queue}{Stack}
\end{question}

\begin{question}{M}
DFS on a graph with V vertices represented by an adjacency matrix takes
\optsshort{O(V + E)}{O(E log V)}{O(V\textsuperscript{2})}{O(V\textsuperscript{3})}
\end{question}

\begin{question}{M}
In a directed graph, the vertex u is discovered at time 1 and finished at time 10, and v is discovered at 3 and finished at 6 during DFS. Which edge type is (u, v) if it exists and v was discovered through another path?
\optsshort{It cannot exist}{Forward edge}{Cross edge}{Back edge}
\end{question}

\begin{question}{M}
Given below are two statements: one is labelled as Assertion A and the other is labelled as Reason R.\\
Assertion A: BFS finds shortest paths in any weighted graph with positive edge weights.\\
Reason R: BFS minimizes the number of edges, not the total weight.
\optslong{Both A and R are true and R is the correct explanation of A}{Both A and R are true but R is NOT the correct explanation of A}{A is true but R is false}{A is false but R is true}
\end{question}

\begin{question}{L}
A graph in which every pair of distinct vertices is joined by an edge is called a
\optsshort{tree}{complete graph}{bipartite graph}{planar graph}
\end{question}

\begin{question}{M}
A network has edges S$\rightarrow$A 3, S$\rightarrow$B 2, A$\rightarrow$B 1, A$\rightarrow$T 2, B$\rightarrow$T 3. The maximum flow from S to T is
\optsshort{4}{7}{6}{5}
\end{question}

\begin{question}{M}
The A* search algorithm finds an optimal path if its heuristic h is
\optslong{always zero only}{admissible (never overestimates the true cost)}{random}{larger than the true cost}
\end{question}

\begin{answers}
\ans{1}{2}{Standard time complexities of graph algorithms.}
\ans{2}{1}{Adjacency lists are preferred for sparse graphs.}
\ans{3}{3}{Kruskal adds g-h 1, c-i 2, f-g 2, a-b 4, c-f 4, c-d 7, a-h 8 (or b-c 8), d-e 9.}
\ans{4}{3}{A-B 1, B-C 1, C-D 2, E-F 3, C-E 5 = 12.}
\ans{5}{3}{A spanning tree has V $-$ 1 edges.}
\ans{6}{4}{Kruskal's time is dominated by sorting the edges.}
\ans{7}{1}{Vertices are extracted in the order s, y, z, t, x.}
\ans{8}{1}{y (5), z (7), t (8), x (9).}
\ans{9}{1}{Bellman-Ford handles negative weights and detects negative cycles.}
\ans{10}{1}{A V-th pass that still improves a distance reveals a negative cycle.}
\ans{11}{4}{The path s $\rightarrow$ y $\rightarrow$ x $\rightarrow$ t $\rightarrow$ z has weight 7 $-$ 3 $-$ 2 $-$ 4 = $-$2.}
\ans{12}{3}{It explores vertices in order of distance from s.}
\ans{13}{2}{DFS goes A $\rightarrow$ B $\rightarrow$ D, backtracks to B $\rightarrow$ E $\rightarrow$ F, and then reaches C from F.}
\ans{14}{4}{In an undirected graph, DFS produces only tree edges and back edges.}
\ans{15}{1}{A must come first and D, E last; B and C can appear in either order.}
\ans{16}{2}{2 and 3 can be swapped, and so can 5 and 6: 2 $\times$ 2 = 4.}
\ans{17}{3}{It can be computed by DFS (decreasing finishing times) or Kahn's algorithm in O(V + E).}
\ans{18}{3}{It runs in O(V + E); Tarjan's algorithm uses a single DFS.}
\ans{19}{3}{The minimum cut \{s, v1, v2, v4\} versus \{v3, t\} has capacity 12 + 7 + 4 = 23.}
\ans{20}{4}{The cut separating T from the rest has capacity A$\rightarrow$T 5 + B$\rightarrow$T 10 = 15, and S$\rightarrow$A$\rightarrow$T 5, S$\rightarrow$A$\rightarrow$B$\rightarrow$T 5, S$\rightarrow$B$\rightarrow$T 5 achieve it.}
\ans{21}{1}{This is the max-flow min-cut theorem.}
\ans{22}{2}{Ford-Fulkerson with arbitrary paths takes O(E |f*|) for integer capacities.}
\ans{23}{2}{Add a source to one side and a sink to the other with unit capacities.}
\ans{24}{1}{1 $\rightarrow$ 2 $\rightarrow$ 3 $\rightarrow$ 4: 3 + 2 + 1 = 6.}
\ans{25}{1}{n(n $-$ 1)/2 = 45.}
\ans{26}{4}{n\textsuperscript{n $-$ 2} = 5\textsuperscript{3} = 125.}
\ans{27}{3}{The Hamiltonian-cycle problem is NP-complete.}
\ans{28}{3}{Bipartiteness can be tested by BFS 2-colouring in O(V + E).}
\ans{29}{2}{Each edge contributes 1 to the degree of each of its two endpoints.}
\ans{30}{1}{Articulation points and bridges can be found in O(V + E) with DFS low-link values.}
\ans{31}{1}{It runs in O(V\textsuperscript{2} log V + VE).}
\ans{32}{1}{Kruskal adds the cheapest edge connecting two different components.}
\ans{33}{4}{It is Floyd-Warshall with Boolean OR and AND.}
\ans{34}{2}{A Fibonacci heap gives O(E + V log V).}
\ans{35}{3}{Each vertex's row of V entries is scanned.}
\ans{36}{2}{v is a descendant of u (its interval is nested inside u's), so a non-tree edge u $\rightarrow$ v is a forward edge.}
\ans{37}{4}{Because of R, BFS is correct only for unweighted graphs (or equal weights); otherwise Dijkstra is needed.}
\ans{38}{2}{K\_n has n(n $-$ 1)/2 edges.}
\ans{39}{4}{The cut \{S\} has capacity 3 + 2 = 5, and S$\rightarrow$A$\rightarrow$T 2, S$\rightarrow$A$\rightarrow$B$\rightarrow$T 1, S$\rightarrow$B$\rightarrow$T 2 achieves 5.}
\ans{40}{2}{With h = 0, A* reduces to Dijkstra's algorithm.}
\end{answers}

\section{Complexity Theory (P, NP, NP-completeness, NP-hardness and reductions)}

\begin{question}{H}
A problem is NP-complete if
\optslong{it is in P}{it is in NP and every problem in NP is polynomial-time reducible to it}{it cannot be solved by any algorithm}{it is NP-hard but not in NP}
\end{question}

\begin{question}{H}
The first problem proved to be NP-complete (by Cook and Levin) is
\optslong{graph colouring}{Boolean satisfiability (SAT)}{subset sum}{the travelling salesman problem}
\end{question}

\begin{question}{H}
Which of the following problems is in P?
\optsshort{Vertex cover}{Hamiltonian cycle}{2-SAT}{3-SAT}
\end{question}

\begin{question}{H}
Which of the following problems are NP-complete?\\
A. Clique  B. Vertex cover  C. Shortest path between two vertices  D. Subset sum  E. Minimum spanning tree\\
Choose the correct answer from the options given below:
\optsshort{B, C and D only}{A, B and D only}{A, B, C, D and E}{A, C and E only}
\end{question}

\begin{question}{H}
Given below are two statements.\\
Statement I: If an NP-complete problem can be solved in polynomial time, then P = NP.\\
Statement II: If P = NP, then every NP-hard problem can be solved in polynomial time.
\optslong{Both Statement I and Statement II are true}{Both Statement I and Statement II are false}{Statement I is true but Statement II is false}{Statement I is false but Statement II is true}
\end{question}

\begin{question}{H}
Problem A is NP-complete, and A is polynomial-time reducible to problem B (A $\leq$p B). Which statement is correct?
\optslong{B is NP-hard}{B is in P}{A is in P}{B is NP-complete in all cases}
\end{question}

\begin{question}{H}
If problem A $\leq$p B and B is in P, then
\optsshort{A is in P}{A is NP-complete}{B is NP-hard}{nothing can be concluded}
\end{question}

\begin{question}{H}
Given below are two statements: one is labelled as Assertion A and the other is labelled as Reason R.\\
Assertion A: The halting problem is NP-hard.\\
Reason R: The halting problem is in NP.
\optslong{Both A and R are true and R is the correct explanation of A}{Both A and R are true but R is NOT the correct explanation of A}{A is true but R is false}{A is false but R is true}
\end{question}

\begin{question}{H}
Match List I with List II.\\
List I (Problem): A. 2-SAT  B. 3-SAT  C. Halting problem  D. Primality testing\\
List II (Class): I. Undecidable  II. In P (AKS algorithm, 2002)  III. NP-complete  IV. In P (linear time via SCCs)
\optsshort{A-IV, B-I, C-III, D-II}{A-IV, B-III, C-II, D-I}{A-III, B-IV, C-I, D-II}{A-IV, B-III, C-I, D-II}
\end{question}

\begin{question}{M}
The 0/1 knapsack problem can be solved in O(nW) time, yet its decision version is NP-complete because
\optslong{O(nW) is polynomial in the input size}{W can be exponential in the number of bits used to represent it, so O(nW) is pseudo-polynomial}{the problem is undecidable}{knapsack is not in NP}
\end{question}

\begin{question}{H}
Given below are two statements.\\
Statement I: A graph G has a clique of size k if and only if its complement has an independent set of size k.\\
Statement II: A set S is a vertex cover of G if and only if V $-$ S is an independent set of G.
\optslong{Both Statement I and Statement II are true}{Both Statement I and Statement II are false}{Statement I is true but Statement II is false}{Statement I is false but Statement II is true}
\end{question}

\begin{question}{M}
Which of the following is NOT known to be NP-complete?
\optslong{Set cover (decision version)}{Hamiltonian path}{Graph 3-colouring}{Integer factorization (decision version)}
\end{question}

\begin{question}{M}
The class NP consists of decision problems
\optslong{that are undecidable}{that cannot be solved in polynomial time}{that are solvable in polynomial time on a deterministic machine}{whose ``yes'' instances have certificates that can be verified in polynomial time}
\end{question}

\begin{question}{M}
Given below are two statements.\\
Statement I: It has been proved that P = NP.\\
Statement II: It has been proved that P $\neq$ NP.
\optslong{Both Statement I and Statement II are true}{Both Statement I and Statement II are false}{Statement I is true but Statement II is false}{Statement I is false but Statement II is true}
\end{question}

\begin{question}{M}
co-NP is the class of problems
\optsshort{that are in P}{that are NP-hard}{whose complements are in NP}{that are not in NP}
\end{question}

\begin{question}{M}
If P = NP, then
\optsshort{nothing follows about co-NP}{NP $\neq$ co-NP}{NP = co-NP}{P $\neq$ co-NP}
\end{question}

\begin{question}{M}
To prove that a new problem X is NP-complete, it is sufficient to show that
\optslong{X is in NP and some known NP-complete problem reduces to X in polynomial time}{X is in P}{X has an exponential algorithm}{X reduces to SAT}
\end{question}

\begin{question}{M}
Match List I with List II.\\
List I (Problem): A. Travelling salesman (decision)  B. Euler circuit  C. Minimum cut in a network  D. Graph isomorphism\\
List II (Status): I. In P by checking that every vertex has even degree (and connectivity)  II. In P via max-flow min-cut  III. NP-complete  IV. In NP, not known to be in P or NP-complete
\optsshort{A-III, B-I, C-II, D-IV}{A-I, B-III, C-II, D-IV}{A-III, B-II, C-I, D-IV}{A-III, B-I, C-IV, D-II}
\end{question}

\begin{question}{M}
A problem is strongly NP-complete if
\optslong{it has a pseudo-polynomial algorithm}{it is in P}{it is undecidable}{it remains NP-complete even when all numbers in its input are bounded by a polynomial in the input length}
\end{question}

\begin{question}{M}
Given below are two statements: one is labelled as Assertion A and the other is labelled as Reason R.\\
Assertion A: 3-SAT is NP-complete.\\
Reason R: SAT can be reduced to 3-SAT in polynomial time by splitting long clauses using new variables.
\optslong{Both A and R are true and R is the correct explanation of A}{Both A and R are true but R is NOT the correct explanation of A}{A is true but R is false}{A is false but R is true}
\end{question}

\begin{question}{M}
The decision version of the travelling salesman problem asks whether
\optslong{there is a tour of total cost at most k}{the graph is connected}{all edges have equal weight}{the graph has an Euler circuit}
\end{question}

\begin{question}{M}
Which of the following problems can be solved in polynomial time?\\
A. Linear programming  B. Integer linear programming  C. Maximum bipartite matching  D. 2-colouring of a graph\\
Choose the correct answer from the options given below:
\optsshort{A, B, C and D}{A, C and D only}{A, B and C only}{B and D only}
\end{question}

\begin{question}{M}
Polynomial-time reductions are transitive. This means that
\optslong{if A $\leq$p B and B $\leq$p C, then A $\leq$p C}{A $\leq$p B implies A = B}{if A $\leq$p B then B $\leq$p A}{every problem reduces to itself in exponential time}
\end{question}

\begin{question}{M}
Given below are two statements.\\
Statement I: Every problem in P is also in co-NP.\\
Statement II: Every NP-complete problem is NP-hard.
\optslong{Both Statement I and Statement II are true}{Both Statement I and Statement II are false}{Statement I is true but Statement II is false}{Statement I is false but Statement II is true}
\end{question}

\begin{question}{L}
The question of whether P = NP was formally introduced in 1971 by
\optsshort{Stephen Cook}{Alan Turing}{Donald Knuth}{Edsger Dijkstra}
\end{question}

\begin{question}{M}
Which of the following is NP-hard but not known to be in NP?
\optslong{3-SAT}{Vertex cover (decision)}{The optimization version of TSP (finding the minimum-cost tour)}{Clique (decision)}
\end{question}

\begin{question}{M}
Karp's 1972 paper showed NP-completeness of how many combinatorial problems?
\optsshort{7}{50}{21}{100}
\end{question}

\begin{answers}
\ans{1}{2}{NP-hard problems need not be in NP; NP-complete problems are NP-hard and in NP.}
\ans{2}{2}{Cook (1971) and Levin independently proved SAT NP-complete.}
\ans{3}{3}{2-SAT is solvable in linear time using strongly connected components of the implication graph.}
\ans{4}{2}{Shortest path and MST have polynomial-time algorithms.}
\ans{5}{3}{Some NP-hard problems, such as the halting problem, are undecidable regardless of whether P = NP.}
\ans{6}{1}{B is at least as hard as A; it is NP-complete only if it is also in NP.}
\ans{7}{1}{A can be solved by transforming it to B in polynomial time and solving B.}
\ans{8}{3}{Every NP problem reduces to the halting problem, but the halting problem is undecidable, so it is not in NP.}
\ans{9}{4}{Primality testing was shown to be in P by Agrawal, Kayal and Saxena.}
\ans{10}{2}{It is weakly NP-complete.}
\ans{11}{1}{These relationships give simple reductions among clique, independent set and vertex cover.}
\ans{12}{4}{Factoring is in NP $\cap$ co-NP and is not believed to be NP-complete.}
\ans{13}{4}{Equivalently, problems solvable in polynomial time by a non-deterministic Turing machine.}
\ans{14}{2}{The question is open; it is one of the Clay Millennium Prize Problems.}
\ans{15}{3}{Tautology (is every assignment satisfying?) is co-NP-complete.}
\ans{16}{3}{P is closed under complement, so P = NP implies NP = co-NP.}
\ans{17}{1}{The reduction must go from the known problem to X.}
\ans{18}{1}{Babai's 2015 result puts graph isomorphism in quasi-polynomial time.}
\ans{19}{4}{Such problems (e.g. 3-partition, TSP) have no pseudo-polynomial algorithm unless P = NP.}
\ans{20}{1}{3-SAT is in NP, and the reduction from SAT shows that it is NP-hard.}
\ans{21}{1}{Optimization versions are NP-hard; decision versions with a bound k are in NP.}
\ans{22}{2}{Integer linear programming is NP-hard.}
\ans{23}{1}{Composing two polynomial-time transformations gives a polynomial-time transformation.}
\ans{24}{1}{P $\subseteq$ NP $\cap$ co-NP, and NP-complete problems are NP-hard by definition.}
\ans{25}{1}{It is one of the Clay Millennium Prize Problems.}
\ans{26}{3}{An optimization problem is not a decision problem, so it is not in NP as stated.}
\ans{27}{3}{Karp's 21 NP-complete problems include clique, vertex cover, Hamiltonian cycle and knapsack.}
\end{answers}

\section{Selected Topics: Number-Theoretic Algorithms, Polynomial Arithmetic, FFT and String Matching}

\begin{question}{H}
gcd(252, 198) computed by Euclid's algorithm is
\optsshort{6}{18}{36}{9}
\end{question}

\begin{question}{H}
The extended Euclidean algorithm applied to 99 and 78 gives gcd = 3 = 99x + 78y with
\optsshort{x = 14, y = $-$11}{x = $-$11, y = 14}{x = $-$3, y = 4}{x = 3, y = $-$4}
\end{question}

\begin{question}{H}
The multiplicative inverse of 7 modulo 26 is
\optsshort{3}{11}{15}{19}
\end{question}

\begin{question}{M}
The multiplicative inverse of 17 modulo 43 is
\optsshort{38}{5}{41}{23}
\end{question}

\begin{question}{H}
The value of 7\textsuperscript{222} mod 11 is
\optsshort{3}{7}{1}{5}
\end{question}

\begin{question}{H}
The value of 2\textsuperscript{100} mod 7 is
\optsshort{2}{4}{1}{6}
\end{question}

\begin{question}{M}
The value of 3\textsuperscript{200} mod 13 is
\optsshort{9}{12}{1}{3}
\end{question}

\begin{question}{H}
The smallest non-negative x satisfying x $\equiv$ 2 (mod 3), x $\equiv$ 3 (mod 5) and x $\equiv$ 2 (mod 7) is
\optsshort{23}{53}{38}{8}
\end{question}

\begin{question}{H}
Euler's totient $\varphi$(100) is
\optsshort{20}{40}{50}{60}
\end{question}

\begin{question}{H}
Given below are two statements.\\
Statement I: Euclid's algorithm takes O(log min(a, b)) division steps.\\
Statement II: The worst case of Euclid's algorithm occurs for consecutive Fibonacci numbers.
\optslong{Both Statement I and Statement II are true}{Both Statement I and Statement II are false}{Statement I is true but Statement II is false}{Statement I is false but Statement II is true}
\end{question}

\begin{question}{M}
The number of division steps that Euclid's algorithm takes to compute gcd(55, 34) is
\optsshort{5}{8}{7}{6}
\end{question}

\begin{question}{H}
Computing a\textsuperscript{b} mod n by repeated squaring requires
\optsshort{O(log b) multiplications}{O(b\textsuperscript{2}) multiplications}{O($\surd$b) multiplications}{O(b) multiplications}
\end{question}

\begin{question}{H}
Given below are two statements: one is labelled as Assertion A and the other is labelled as Reason R.\\
Assertion A: Fermat's little theorem states that a\textsuperscript{p $-$ 1} $\equiv$ 1 (mod p) for a prime p not dividing a.\\
Reason R: The multiplicative group of integers modulo a prime p has order p $-$ 1.
\optslong{Both A and R are true and R is the correct explanation of A}{Both A and R are true but R is NOT the correct explanation of A}{A is true but R is false}{A is false but R is true}
\end{question}

\begin{question}{M}
The multiplicative inverse of a modulo n exists if and only if
\optsshort{a is odd}{n is prime}{a < n}{gcd(a, n) = 1}
\end{question}

\begin{question}{H}
Evaluating a polynomial of degree n at a point using Horner's rule requires
\optslong{n(n + 1)/2 multiplications}{2n multiplications and n additions}{n multiplications and n additions}{n\textsuperscript{2} multiplications}
\end{question}

\begin{question}{H}
Multiplying two polynomials of degree n using the fast Fourier transform takes
\optsshort{O(n log\textsuperscript{2} n)}{O(n\textsuperscript{2})}{O(n\textsuperscript{1}.58)}{O(n log n)}
\end{question}

\begin{question}{H}
Given below are two statements.\\
Statement I: The FFT evaluates a polynomial at the complex nth roots of unity in O(n log n) time.\\
Statement II: The naive evaluation of a degree-(n $-$ 1) polynomial at n points takes O(n log n) time.
\optslong{Both Statement I and Statement II are true}{Both Statement I and Statement II are false}{Statement I is true but Statement II is false}{Statement I is false but Statement II is true}
\end{question}

\begin{question}{M}
The FFT algorithm is based on which technique?
\optsshort{Greedy}{Divide and conquer}{Backtracking}{Dynamic programming}
\end{question}

\begin{question}{M}
The FFT requires the number of evaluation points n to be
\optslong{a power of 2 (in the basic radix-2 version)}{odd}{prime}{a multiple of 10}
\end{question}

\begin{question}{H}
Match List I with List II.\\
List I (String-matching algorithm): A. Naive  B. Rabin-Karp  C. Knuth-Morris-Pratt  D. Finite-automaton matcher\\
List II (Matching time, text length n, pattern length m): I. O((n $-$ m + 1)m) worst case, O(n + m) expected  II. O(n) after O(m |$\Sigma$|) preprocessing  III. O((n $-$ m + 1)m)  IV. O(n) after O(m) preprocessing
\optsshort{A-III, B-I, C-IV, D-II}{A-III, B-IV, C-I, D-II}{A-I, B-III, C-IV, D-II}{A-III, B-I, C-II, D-IV}
\end{question}

\begin{question}{H}
The KMP prefix function $\pi$ for the pattern ``ababaca'' is
\optsshort{0 1 2 3 0 1 0}{0 0 1 2 0 0 1}{0 0 0 1 2 3 0}{0 0 1 2 3 0 1}
\end{question}

\begin{question}{M}
The KMP prefix function $\pi$ for the pattern ``aabaaab'' is
\optsshort{0 0 1 1 2 2 3}{0 1 0 1 2 3 3}{0 1 0 1 2 2 3}{0 1 0 0 1 2 3}
\end{question}

\begin{question}{M}
The KMP prefix function $\pi$ for the pattern ``abcdabd'' is
\optsshort{0 0 0 0 1 2 0}{0 1 2 3 0 0 0}{0 0 0 0 1 2 3}{0 0 0 1 2 0 0}
\end{question}

\begin{question}{H}
In the Rabin-Karp example with text 2359023141526739921, pattern 31415 and modulus 13, the number of spurious hits is
\optsshort{0}{1}{2}{3}
\end{question}

\begin{question}{H}
Given below are two statements: one is labelled as Assertion A and the other is labelled as Reason R.\\
Assertion A: Rabin-Karp must verify each hash match character by character.\\
Reason R: Different strings can have the same hash value modulo q.
\optslong{Both A and R are true and R is the correct explanation of A}{Both A and R are true but R is NOT the correct explanation of A}{A is true but R is false}{A is false but R is true}
\end{question}

\begin{question}{M}
The naive string-matching algorithm, comparing characters left to right and stopping at the first mismatch, makes how many character comparisons to search for ``aaab'' in ``aaaaaaaaab''?
\optsshort{10}{7}{28}{40}
\end{question}

\begin{question}{M}
Given below are two statements.\\
Statement I: The Boyer-Moore algorithm scans the pattern from right to left.\\
Statement II: Boyer-Moore can skip parts of the text, so it is often sublinear in practice.
\optslong{Both Statement I and Statement II are true}{Both Statement I and Statement II are false}{Statement I is true but Statement II is false}{Statement I is false but Statement II is true}
\end{question}

\begin{question}{M}
The KMP algorithm never
\optsshort{moves backwards in the text}{preprocesses the pattern}{matches in linear time}{uses the prefix function}
\end{question}

\begin{question}{M}
The Rabin-Karp algorithm is especially useful for
\optslong{compressing text}{sorting strings}{searching for many patterns at once, and for 2-D pattern matching}{single-character search}
\end{question}

\begin{question}{M}
Given below are two statements.\\
Statement I: The Miller-Rabin primality test is a deterministic polynomial-time algorithm that never errs.\\
Statement II: Modular exponentiation is used in the RSA cryptosystem.
\optslong{Both Statement I and Statement II are true}{Both Statement I and Statement II are false}{Statement I is true but Statement II is false}{Statement I is false but Statement II is true}
\end{question}

\begin{question}{M}
The value of 5\textsuperscript{117} mod 19 is
\optsshort{1}{7}{11}{5}
\end{question}

\begin{question}{M}
The number of operations needed to multiply two n-digit integers with the schoolbook method is
\optsshort{O(n\textsuperscript{2})}{O(n)}{O(2\textsuperscript{n})}{O(n log n)}
\end{question}

\begin{question}{M}
A suffix tree for a text of length n can be built in
\optsshort{O(n) time}{O(n\textsuperscript{2}) time only}{O(2\textsuperscript{n}) time}{O(n\textsuperscript{3}) time}
\end{question}

\begin{question}{L}
The greatest common divisor of two consecutive integers is always
\optsshort{the smaller integer}{1}{2}{0}
\end{question}

\begin{question}{M}
Given below are two statements: one is labelled as Assertion A and the other is labelled as Reason R.\\
Assertion A: The inverse DFT can be computed with the same FFT algorithm.\\
Reason R: The inverse DFT requires $\Theta$(n\textsuperscript{2}) time.
\optslong{Both A and R are true and R is the correct explanation of A}{Both A and R are true but R is NOT the correct explanation of A}{A is true but R is false}{A is false but R is true}
\end{question}

\begin{answers}
\ans{1}{2}{252 = 1 $\times$ 198 + 54; 198 = 3 $\times$ 54 + 36; 54 = 1 $\times$ 36 + 18; 36 = 2 $\times$ 18.}
\ans{2}{2}{99 $\times$ ($-$11) + 78 $\times$ 14 = $-$1089 + 1092 = 3.}
\ans{3}{3}{7 $\times$ 15 = 105 = 4 $\times$ 26 + 1.}
\ans{4}{1}{17 $\times$ 38 = 646 = 15 $\times$ 43 + 1.}
\ans{5}{4}{By Fermat, 7\textsuperscript{10} $\equiv$ 1 (mod 11); 7\textsuperscript{222} = (7\textsuperscript{10})\textsuperscript{22} $\times$ 7\textsuperscript{2} $\equiv$ 49 $\equiv$ 5.}
\ans{6}{1}{2\textsuperscript{3} $\equiv$ 1 (mod 7); 100 = 3 $\times$ 33 + 1, so 2\textsuperscript{100} $\equiv$ 2.}
\ans{7}{1}{3\textsuperscript{3} = 27 $\equiv$ 1 (mod 13); 200 = 3 $\times$ 66 + 2, so 3\textsuperscript{200} $\equiv$ 9.}
\ans{8}{1}{By the Chinese remainder theorem the solution is unique modulo 105; 23 satisfies all three.}
\ans{9}{2}{$\varphi$(100) = 100 $\times$ (1 $-$ 1/2) $\times$ (1 $-$ 1/5) = 40.}
\ans{10}{1}{This is Lam\'e's theorem.}
\ans{11}{2}{55 and 34 are consecutive Fibonacci numbers, giving the maximum number of steps for their size: 8.}
\ans{12}{1}{The exponent is processed bit by bit.}
\ans{13}{1}{Lagrange's theorem then gives a\textsuperscript{p $-$ 1} $\equiv$ 1.}
\ans{14}{4}{It is computed with the extended Euclidean algorithm.}
\ans{15}{3}{a0 + x(a1 + x(a2 + ... + x an)).}
\ans{16}{4}{Evaluate at 2n roots of unity with FFT, multiply pointwise, and interpolate with inverse FFT.}
\ans{17}{3}{Naive evaluation at n points with Horner's rule takes O(n\textsuperscript{2}).}
\ans{18}{2}{A(x) = A\_even(x\textsuperscript{2}) + x A\_odd(x\textsuperscript{2}) splits the problem into two halves.}
\ans{19}{1}{The polynomial is padded with zero coefficients if needed.}
\ans{20}{1}{Standard complexities from CLRS.}
\ans{21}{4}{$\pi$[q] is the length of the longest proper prefix of P[1..q] that is also its suffix.}
\ans{22}{3}{For example, ``aabaa'' has border ``aa'' (2), and ``aabaaab'' has border ``aab'' (3).}
\ans{23}{1}{``abcdab'' has border ``ab''; ``abcdabd'' has no non-empty border.}
\ans{24}{2}{The valid match is at shift 6; the window 67399 at shift 12 also has hash 7 but does not match.}
\ans{25}{1}{Such false matches are called spurious hits.}
\ans{26}{3}{Each of the 7 shifts compares 4 characters.}
\ans{27}{1}{It uses the bad-character and good-suffix heuristics.}
\ans{28}{1}{The text pointer only advances; the pattern pointer falls back using $\pi$.}
\ans{29}{3}{The hashes of all patterns can be stored in a set.}
\ans{30}{4}{Miller-Rabin is a randomized (Monte Carlo) test with a small probability of error for composites.}
\ans{31}{1}{5 has order 9 modulo 19 (5\textsuperscript{9} $\equiv$ 1), and 117 = 9 $\times$ 13.}
\ans{32}{1}{Karatsuba reduces it to O(n\textsuperscript{1}.585), and FFT-based methods to about O(n log n).}
\ans{33}{1}{Ukkonen's and McCreight's algorithms are linear for a constant alphabet.}
\ans{34}{2}{Any common divisor would also divide their difference, 1.}
\ans{35}{3}{The inverse DFT uses conjugate roots of unity and division by n, so it also takes O(n log n).}
\end{answers}

\section{Advanced Algorithms: Parallel, Approximation and Randomized Algorithms}

\begin{question}{H}
Match List I with List II.\\
List I (PRAM model): A. EREW  B. CREW  C. CRCW\\
List II (Memory access rule): I. Several processors may read and write the same cell at the same time, with a rule to resolve write conflicts  II. No two processors may read or write the same cell at the same time  III. Several processors may read the same cell at the same time, but writes must be exclusive
\optsshort{A-II, B-I, C-III}{A-I, B-III, C-II}{A-III, B-II, C-I}{A-II, B-III, C-I}
\end{question}

\begin{question}{H}
The sum of n numbers can be computed on a PRAM with n/2 processors in
\optsshort{O(1) time}{O(log n) time}{O(n log n) time}{O(n) time}
\end{question}

\begin{question}{H}
The speedup and efficiency of a parallel algorithm using p processors are defined as
\optsshort{T1 $\times$ Tp and T1/p}{T1/Tp and (T1/Tp)/p}{Tp/T1 and p/Tp}{p/T1 and T1/p}
\end{question}

\begin{question}{M}
A parallel algorithm takes 100 s on 1 processor and 20 s on 8 processors. Its speedup and efficiency are
\optsshort{8 and 100\%}{5 and 62.5\%}{20 and 250\%}{5 and 80\%}
\end{question}

\begin{question}{M}
Brent's theorem states that a computation with total work W and depth (span) D can be executed on p processors in time
\optsshort{W + D $\times$ p}{W/p + D}{W $\times$ p}{D/p}
\end{question}

\begin{question}{H}
Bitonic sort on n elements, as a sorting network, has depth
\optsshort{O(n log n)}{O(log\textsuperscript{2} n)}{O(log n)}{O(n)}
\end{question}

\begin{question}{M}
Odd-even transposition sort on a linear array of n processors sorts n elements in
\optsshort{$\surd$n phases}{n phases}{n\textsuperscript{2} phases}{log n phases}
\end{question}

\begin{question}{M}
The parallel prefix-sum (scan) of n elements can be computed in
\optslong{O(1) time}{O(log n) time with O(n) work}{O(n) time}{O(n log n) time with O(n\textsuperscript{2}) work}
\end{question}

\begin{question}{M}
Given below are two statements.\\
Statement I: A parallel algorithm is work-optimal if its total work matches the best sequential running time.\\
Statement II: Adding more processors always reduces the running time of a parallel algorithm.
\optslong{Both Statement I and Statement II are true}{Both Statement I and Statement II are false}{Statement I is true but Statement II is false}{Statement I is false but Statement II is true}
\end{question}

\begin{question}{M}
On a CREW PRAM with p processors, searching a sorted array of n elements takes
\optsshort{O(log n / log(p + 1))}{O(n/p)}{O(log n)}{O(1)}
\end{question}

\begin{question}{H}
The approximation ratio of the simple maximal-matching algorithm for the minimum vertex-cover problem is
\optsshort{2}{1}{1.5}{log n}
\end{question}

\begin{question}{H}
For the travelling salesman problem with the triangle inequality, the MST-based (double-tree) algorithm gives an approximation ratio of
\optsshort{1.5}{log n}{1}{2}
\end{question}

\begin{question}{H}
Given below are two statements: one is labelled as Assertion A and the other is labelled as Reason R.\\
Assertion A: The general travelling salesman problem (without the triangle inequality) has no polynomial-time constant-factor approximation unless P = NP.\\
Reason R: Such an approximation could be used to decide the Hamiltonian-cycle problem in polynomial time.
\optslong{Both A and R are true and R is the correct explanation of A}{Both A and R are true but R is NOT the correct explanation of A}{A is true but R is false}{A is false but R is true}
\end{question}

\begin{question}{M}
The greedy algorithm for set cover has an approximation ratio of
\optsshort{H(n) $\approx$ ln n}{2}{1.5}{n}
\end{question}

\begin{question}{M}
Match List I with List II.\\
List I (Term): A. PTAS  B. FPTAS  C. Approximation ratio $\rho$(n)\\
List II (Meaning): I. For every $\varepsilon$ > 0, a (1 + $\varepsilon$)-approximation in time polynomial in n and 1/$\varepsilon$  II. For every $\varepsilon$ > 0, a (1 + $\varepsilon$)-approximation in time polynomial in n for each fixed $\varepsilon$  III. Bound on the ratio between the cost of the produced solution and the optimal cost
\optsshort{A-II, B-III, C-I}{A-III, B-I, C-II}{A-II, B-I, C-III}{A-I, B-II, C-III}
\end{question}

\begin{question}{M}
The first-fit decreasing heuristic for bin packing uses at most about
\optsshort{11/9 OPT + constant bins}{2 OPT bins}{OPT bins}{log n $\times$ OPT bins}
\end{question}

\begin{question}{H}
Match List I with List II.\\
List I (Randomized algorithm type): A. Las Vegas  B. Monte Carlo\\
List II (Property): I. Always gives the correct answer; running time is random  II. Running time is bounded; the answer may be wrong with small probability
\optsshort{A-I, B-II}{A-II, B-I}{Both I}{Both II}
\end{question}

\begin{question}{H}
The expected running time of randomized quick sort on any input of n distinct elements is
\optsshort{O(n)}{O(n log n)}{O(log n)}{O(n\textsuperscript{2})}
\end{question}

\begin{question}{M}
Karger's randomized contraction algorithm finds a minimum cut of a graph with n vertices in one run with probability at least
\optsshort{2/(n(n $-$ 1))}{1/n}{1/2}{1 $-$ 1/n}
\end{question}

\begin{question}{M}
Given below are two statements.\\
Statement I: Freivalds' algorithm checks whether AB = C for n $\times$ n matrices in O(n\textsuperscript{2}) time per trial.\\
Statement II: Freivalds' algorithm can wrongly report that AB = C, with probability at most 1/2 per trial when AB $\neq$ C.
\optslong{Both Statement I and Statement II are true}{Both Statement I and Statement II are false}{Statement I is true but Statement II is false}{Statement I is false but Statement II is true}
\end{question}

\begin{question}{M}
In the hiring problem with n candidates in random order, the expected number of times a new best candidate is hired is
\optsshort{1}{n}{n/2}{H(n) $\approx$ ln n}
\end{question}

\begin{question}{M}
Match List I with List II.\\
List I (Complexity class): A. RP  B. ZPP  C. BPP\\
List II (Definition): I. Bounded two-sided error, probability of error at most 1/3  II. One-sided error: never accepts a ``no'' instance and accepts a ``yes'' instance with probability at least 1/2  III. Zero error in expected polynomial time (Las Vegas)
\optsshort{A-II, B-III, C-I}{A-II, B-I, C-III}{A-I, B-III, C-II}{A-III, B-II, C-I}
\end{question}

\begin{question}{M}
The expected search time in a skip list with n elements is
\optsshort{O(1)}{O($\surd$n)}{O(log n)}{O(n)}
\end{question}

\begin{question}{M}
Given below are two statements: one is labelled as Assertion A and the other is labelled as Reason R.\\
Assertion A: Randomized algorithms can avoid worst-case inputs that are designed against a deterministic algorithm.\\
Reason R: The behaviour of a randomized algorithm depends on random choices made at run time, not only on the input.
\optslong{Both A and R are true and R is the correct explanation of A}{Both A and R are true but R is NOT the correct explanation of A}{A is true but R is false}{A is false but R is true}
\end{question}

\begin{question}{M}
The Fisher-Yates algorithm generates a uniformly random permutation of n elements in
\optsshort{O(n) time}{O(n log n) time}{O(n\textsuperscript{2}) time}{O(2\textsuperscript{n}) time}
\end{question}

\begin{question}{M}
Randomized selection (quickselect with random pivots) finds the kth smallest of n elements in
\optsshort{O(log n) expected time}{O(1) time}{O(n) expected time}{O(n log n) worst-case time}
\end{question}

\begin{question}{M}
Given below are two statements.\\
Statement I: An approximation algorithm always finds the optimal solution.\\
Statement II: An approximation algorithm runs in polynomial time and returns a solution within a guaranteed factor of the optimum.
\optslong{Both Statement I and Statement II are true}{Both Statement I and Statement II are false}{Statement I is true but Statement II is false}{Statement I is false but Statement II is true}
\end{question}

\begin{question}{M}
For the MAX-CUT problem, the Goemans-Williamson algorithm based on semidefinite programming achieves an approximation ratio of about
\optsshort{0.5}{0.878}{1}{0.95}
\end{question}

\begin{question}{L}
In parallel computing, the span (critical path length) of a computation is
\optslong{the number of processors}{the length of the longest chain of dependent operations}{the memory used}{the total number of operations}
\end{question}

\begin{question}{M}
A Bloom filter with m bits, k hash functions and n inserted elements has a false-positive probability of about
\optsshort{e\textsuperscript{$-$k}}{(1 $-$ e\textsuperscript{$-$kn/m})\textsuperscript{k}}{1/m}{kn/m}
\end{question}

\begin{question}{M}
Given below are two statements.\\
Statement I: A Monte Carlo algorithm for a decision problem can be converted into one with smaller error probability by repeating it independently.\\
Statement II: A Las Vegas algorithm may return an incorrect answer.
\optslong{Both Statement I and Statement II are true}{Both Statement I and Statement II are false}{Statement I is true but Statement II is false}{Statement I is false but Statement II is true}
\end{question}

\begin{answers}
\ans{1}{4}{CRCW variants include common, arbitrary and priority write rules.}
\ans{2}{2}{Numbers are added in pairs in a balanced binary tree of depth log2 n.}
\ans{3}{2}{T1 is the best sequential time and Tp the parallel time on p processors.}
\ans{4}{2}{Speedup = 100/20 = 5; efficiency = 5/8 = 62.5\%.}
\ans{5}{2}{More precisely, T\_p $\leq$ W/p + D.}
\ans{6}{2}{It performs O(n log\textsuperscript{2} n) comparisons in total.}
\ans{7}{2}{Odd and even pairs are compared alternately.}
\ans{8}{2}{Up-sweep and down-sweep phases each take log n steps.}
\ans{9}{3}{The span and communication overhead limit the speedup (see Amdahl's law).}
\ans{10}{1}{Each step splits the array into p + 1 parts, so p = 1 gives binary search.}
\ans{11}{1}{Taking both endpoints of every edge in a maximal matching gives a cover at most twice the optimum.}
\ans{12}{4}{Christofides' algorithm improves this to 1.5.}
\ans{13}{1}{Give non-graph edges a huge weight; a good approximation would reveal whether a Hamiltonian cycle exists.}
\ans{14}{1}{It repeatedly picks the set covering the most uncovered elements.}
\ans{15}{3}{0/1 knapsack has an FPTAS.}
\ans{16}{1}{Simple first fit uses at most about 1.7 OPT bins.}
\ans{17}{1}{Randomized quick sort is Las Vegas; Miller-Rabin is Monte Carlo.}
\ans{18}{2}{The expected number of comparisons is about 2n ln n $\approx$ 1.39 n log2 n.}
\ans{19}{1}{Repeating it about n\textsuperscript{2} ln n times makes failure unlikely.}
\ans{20}{1}{Repeating k trials reduces the error probability to 2\textsuperscript{$-$k}.}
\ans{21}{4}{Candidate i is the best so far with probability 1/i.}
\ans{22}{1}{ZPP = RP $\cap$ co-RP.}
\ans{23}{3}{Each element is promoted to the next level with probability 1/2.}
\ans{24}{1}{An adversary cannot predict the random pivot of randomized quick sort.}
\ans{25}{1}{For i from n down to 2, swap A[i] with A[random(1, i)].}
\ans{26}{3}{Its worst case is O(n\textsuperscript{2}), unlike median of medians.}
\ans{27}{4}{Approximation algorithms trade optimality for efficiency.}
\ans{28}{2}{A random cut already achieves 0.5 in expectation.}
\ans{29}{2}{It is the running time with infinitely many processors.}
\ans{30}{2}{The best k is about (m/n) ln 2.}
\ans{31}{3}{A Las Vegas algorithm is always correct; only its running time varies.}
\end{answers}
